\documentclass[10pt,a4paper]{article}
\usepackage{amsmath,amssymb,bm,graphicx,makeidx,subfigure}
\usepackage[italian,english]{babel}
\usepackage[center,small]{caption}[2007/01/07]
\usepackage{fancyhdr}
\usepackage{color}
\usepackage{hyperref}
\hypersetup{
    colorlinks=true, %set true if you want colored links
    linktoc=all,     %set to all if you want both sections and subsections linked
    linkcolor=blue,  %choose some color if you want links to stand out
}

\oddsidemargin = 12pt
\topmargin = 0pt
\textwidth = 440pt
\textheight = 650pt

\makeindex

\begin{document}

\title{Matching conditions}

\tableofcontents
\newpage

\section{PDF matching conditions}

If the (Zero Mass) Variable Flavour Number Scheme (ZM$-$VFNS) at NNLO
is considered, matching conditions for the PDFs and the coupling
constant at the heavy-quarks thresholds ($m_c^2$, $m_b^2$, and $m_t^2$)
must be implemented. This is due to the fact that we are working with
an ``effective theory'' where before a certain threshold, say $m_h^2$,
the heavy-quark flavour $h$ is treated as infinitely massive, while
after the crossing of the same threshold the same flavour is treated
as massless. This results in discontinuities of PDFs and coupling
constant in correspondence of the thresholds. Such discontinuities can
be evaluated in perturbation theory~\cite{Buza:1996wv}.

The discontinuity of the PDF $l$ (in Mellin space) of a light
quark(\footnote{Note that the light quarks run from $1$ to $n_f$.})
just beyond the threshold $m_h^2$ ($=m_c^2,m_b^2,m_t^2$), where the
effective flavour number goes from $n_f$ to $n_f+1$, is given as a
function of the same PDF just before the threshold by the following
relation \cite{Vogt:2004ns}:
\begin{equation}
l^{(n_f+1)}(N,m_h^2)=[1+a_s^2(m_h^2)A_{qq,h}^{N\!S,(2)}(N)]l^{(n_f)}(N,m_h^2)\,.
\label{eq:lightqmc}
\end{equation}
with $l=u,\overline{u},d,\overline{d},\dots$, while the gluon
distribution function is given by:
\begin{equation}
\displaystyle
g^{(n_f+1)}(N,m_h^2)=[1+a_s^2(m_h^2)A_{gg,h}^{S,(2)}(N)]g^{(n_f)}(N,m_h^2)+a_s^2(m_h^2)A^{S,(2)}_{gq,h}(N)\Sigma^{(n_f)}(N,m^2_h)
\label{gluon}
\end{equation}
and, in the end, the sum of heavy quark $h$ and its anti-quark
$\overline{h}$, which are going to be produced after the threshold
$m_h^2$, is:
\begin{equation}
(h^{(n_f+1)}+\overline{h}^{(n_f+1)})(N,m_h^2)=a_s^2(m_h^2)[\tilde{A}^{S,(2)}_{hq}(N)\Sigma^{(n_f)}(N,m_h^2)+\tilde{A}^{S,(2)}_{hg}(N)g^{(n_f)}(N,m_h^2)]\,.
\label{eq:heavyqmc}
\end{equation}
Of course, we have $h=\overline{h}$. 

Now, since:
\begin{equation}
\Sigma^{(n_f+1)}=\sum_{l=1}^{n_f}(l^{(n_f+1)}+\overline{l}^{(n_f+1)})+(h^{(n_f+1)}+\overline{h}^{(n_f+1)})
\end{equation}
we find that the matching condition for the singlet is:
\begin{equation}
\begin{array}{c}
\displaystyle\Sigma^{(n_f+1)}(N,m_h^2)=[1+a_s^2(m_h^2)A_{qq,h}^{N\!S,(2)}(N)]\Sigma^{(n_f)}(N,m_h^2)+\\
\\
\displaystyle a_s^2(m_h^2)[\tilde{A}^{S,(2)}_{hq}(N)\Sigma^{(n_f)}(N,m_h^2)+\tilde{A}^{S,(2)}_{hg}(N)g^{(n_f)}(N,m_h^2)]
\end{array}
\label{singlet}
\end{equation}
so that, from eqs. (\ref{gluon}) and (\ref{singlet}):
\begin{equation}
\begin{array}{c}
\displaystyle {\Sigma^{(n_f+1)} \choose g^{(n_f+1)}}=\begin{pmatrix}1+a_s^2[A_{qq,h}^{N\!S,(2)}+\tilde{A}^{S,(2)}_{hq}] & a_s^2\tilde{A}^{S,(2)}_{hg}\\
a_s^2A^{S,(2)}_{gq,h} & 1+a_s^2A_{gg,h}^{S,(2)}\end{pmatrix}{\Sigma^{(n_f)} \choose g^{(n_f)}}\\
\\
\displaystyle =\left[\begin{pmatrix} 1 & 0 \\ 0 & 1\end{pmatrix}+a_s^2A_{qq,h}^{N\!S,(2)}\begin{pmatrix} 1 & 0 \\ 0 & 0\end{pmatrix}+a_s^2\begin{pmatrix} \tilde{A}^{S,(2)}_{hq} & \tilde{A}^{S,(2)}_{hg} \\A^{S,(2)}_{gq,h} & A_{gg,h}^{S,(2)}\end{pmatrix}\right]{\Sigma^{(n_f)} \choose g^{(n_f)}}
\end{array}
\label{couple1}
\end{equation}
where we have omitted all the dependencies. So we have obtained the
matching conditions for the singlet and the gluon distribution
functions.

Now we consider the other distributions. The valence distribution $V$
for $n_f+1$ active (light) flavours just beyond the threshold $m_h^2$
is defined as:
\begin{equation}
V^{(n_f+1)}=\sum_{l=1}^{n_f}(l^{(n_f+1)}-\overline{l}^{(n_f+1)})+(h^{(n_f+1)}-\overline{h}^{(n_f+1)})
\end{equation}
but, since $h=\overline{h}$, the last term vanish and we are left with:
\begin{equation}
\begin{array}{c}
\displaystyle V^{(n_f+1)}=\sum_{l=1}^{n_f}(l^{(n_f+1)}-\overline{l}^{(n_f+1)})=\\
\\
\displaystyle [1+a_s^2A_{qq,h}^{N\!S,(2)}]\sum_{l=1}^{n_f}(l^{(n_f)}-\overline{l}^{(n_f)})=[1+a_s^2A_{qq,h}^{N\!S,(2)}]V^{(n_f)}\,.
\end{array}
\label{valence}
\end{equation}
So, this is the matching condition for the valence distribution $V$.

Now we consider the valence distribution $V_3$ and $V_8$ which are
both composed only by light quarks, namely:
\begin{equation}
V_3=(u-\overline{u})-(d-\overline{d})\quad\mbox{and}\quad V_8=(u-\overline{u})+(d-\overline{d})-2(s-\overline{s})
\end{equation}
so that, in these cases, the matching conditions work as in the case
of $V$, i.e.:
\begin{equation}
V_{3,8}^{(n_f+1)}=[1+a_s^2A_{qq,h}^{N\!S,(2)}]V^{(n_f)}_{3,8}\,.
\label{v38}
\end{equation}

The same holds for $T_{3}$ and $T_8$, which are defined as:
\begin{equation}
T_3=(u+\overline{u})-(d+\overline{d})\quad\mbox{and}\quad V_8=(u+\overline{u})+(d+\overline{d})-2(s+\overline{s})
\end{equation}
so:
\begin{equation}
T_{3,8}^{(n_f+1)}=[1+a_s^2A_{qq,h}^{N\!S,(2)}]T^{(n_f)}_{3,8}\,.
\label{t38}
\end{equation}

The remaining valence distribution $V_{15}$, $V_{24}$ and $V_{35}$ are
defined as:
\begin{equation}
\begin{array}{l}
V_{15}=(u-\overline{u})+(d-\overline{d})+(s-\overline{s})-3(c-\overline{c})\\
V_{24}=(u-\overline{u})+(d-\overline{d})+(s-\overline{s})+(c-\overline{c})-4(b-\overline{b})\\
V_{35}=(u-\overline{u})+(d-\overline{d})+(s-\overline{s})+(c-\overline{c})+(b-\overline{b})-5(t-\overline{t})
\end{array}
\end{equation}
and since in each one of them the heavy quarks appear always as
difference between quark and anti-quark, they cancel exactly. For
example, at the $m_b^2$ threshold, $V_{15}$ is entirely composed by
light quarks so there is no problem, while $V_{24}$ and $V_{35}$ have
also a $b$-quark contribution, given by $-4(b-\overline{b})$ and
$(b-\overline{b})$ respectively (of course, the $t(\overline{t})$
distribution is zero). Anyway, this terms give no matching condition
since the $b$ contribution is exactly equal to the $\overline{b}$
contribution, so that they cancel. So:
\begin{equation}
V_{15,24,35}^{(n_f+1)}=[1+a_s^2A_{qq,h}^{N\!S,(2)}]V^{(n_f)}_{15,24,35}\,.
\label{v152435}
\end{equation}

In the end, to deal with $T_{15}$, $T_{24}$ and $T_{35}$, we have to
specify the threshold. Indeed, in these cases the heavy quark
contribution does not cancel. Their definition is:
\begin{equation}
\begin{array}{l}
T_{15}=(u+\overline{u})+(d+\overline{d})+(s+\overline{s})-3(c+\overline{c})\\
T_{24}=(u+\overline{u})+(d+\overline{d})+(s+\overline{s})+(c+\overline{c})-4(b+\overline{b})\\
T_{35}=(u+\overline{u})+(d+\overline{d})+(s+\overline{s})+(c+\overline{c})+(b+\overline{b})-5(t+\overline{t})\,.
\end{array}
\end{equation}
Just before the threshold $m_c^2$ we have only 3 active light flavours
($u$, $d$ and $s$), while just beyond $m_c^2$ we have 4 active
flavours and among them the flavour $c$ is considered to be heavy. Of
course, we have no $b$ and $t$ contribution (so $T_{24}$ and $T_{35}$
are equal). So the matching conditions are:
\begin{equation}
\begin{array}{c}
\displaystyle T_{15}^{(4)}=[1+a_s^2A_{qq,c}^{N\!S,(2)}]\underbrace{\sum_{l=u,d,s}(l^{(3)}+\overline{l}^{(3)})}_{\Sigma^{(3)}}-3a_s^2[\tilde{A}^{S,(2)}_{cq}\Sigma^{(3)}+\tilde{A}^{S,(2)}_{cg}g^{(3)}]=\\
\\
\displaystyle \begin{pmatrix} 1+a_s^2[A_{qq,c}^{N\!S,(2)}-3\tilde{A}^{S,(2)}_{cq}] & -3a_s^2\tilde{A}^{S,(2)}_{cg}\end{pmatrix}{\Sigma^{(3)} \choose g^{(3)}}
\end{array}
\end{equation}
while:
\begin{equation}
T_{24,35}^{(4)}=\begin{pmatrix} 1+a_s^2[A_{qq,c}^{N\!S,(2)}+\tilde{A}^{S,(2)}_{cq}] & a_s^2\tilde{A}^{S,(2)}_{cg}\end{pmatrix}{\Sigma^{(3)} \choose g^{(3)}}\,.
\end{equation}
We can put the above relation in a matricial form:
\begin{equation}
\begin{pmatrix} T_{15}^{(4)} \\ T_{24}^{(4)} \\ T_{35}^{(4)} \end{pmatrix} = \begin{pmatrix}  1+a_s^2[A_{qq,c}^{N\!S,(2)}-3\tilde{A}^{S,(2)}_{cq}] & -3a_s^2\tilde{A}^{S,(2)}_{cg}\\  
1+a_s^2[A_{qq,c}^{N\!S,(2)}+\tilde{A}^{S,(2)}_{cq}] & a_s^2\tilde{A}^{S,(2)}_{cg} \\
1+a_s^2[A_{qq,c}^{N\!S,(2)}+\tilde{A}^{S,(2)}_{cq}] &
a_s^2\tilde{A}^{S,(2)}_{cg} \end{pmatrix}{\Sigma^{(3)} \choose
g^{(3)}}\,.
\label{pippo1}
\end{equation}
But now it is easy to generalize. At $m_b^2$, $T_{15}$ does not
contain, so:
\begin{equation}
T_{15}^{(5)}=[1+a_s^2A_{qq,b}^{N\!S,(2)}]T_{15}^{(4)}
\label{pippo2}
\end{equation}
while:
\begin{equation}
\begin{pmatrix} T_{24}^{(5)} \\ T_{35}^{(5)} \end{pmatrix} = \begin{pmatrix}  1+a_s^2[A_{qq,b}^{N\!S,(2)}-4\tilde{A}^{S,(2)}_{bq}] & -4a_s^2\tilde{A}^{S,(2)}_{bg}\\  
1+a_s^2[A_{qq,b}^{N\!S,(2)}+\tilde{A}^{S,(2)}_{bq}] &
a_s^2\tilde{A}^{S,(2)}_{bg} \end{pmatrix}{\Sigma^{(4)} \choose
g^{(4)}}\,.
\label{pippo3}
\end{equation}
Finally, at $m_t^2$ we have:
\begin{equation}
\begin{array}{l}
\displaystyle T_{15}^{(6)}=[1+a_s^2A_{qq,t}^{N\!S,(2)}]T_{15}^{(5)}\\
\\
\displaystyle T_{24}^{(6)}=[1+a_s^2A_{qq,t}^{N\!S,(2)}]T_{25}^{(5)}
\end{array}
\label{pippo4}
\end{equation}
and:
\begin{equation}
T_{35}^{(6)} = \begin{pmatrix}
  1+a_s^2[A_{qq,t}^{N\!S,(2)}-5\tilde{A}^{S,(2)}_{tq}] &
  -5a_s^2\tilde{A}^{S,(2)}_{tg}\end{pmatrix}{\Sigma^{(5)} \choose
  g^{(5)}}\,.
\label{pippo5}
\end{equation}

An explicit calculation for the coefficients $A^{(2)}$ in the
$x$-space can be found in [hep-ph/9612398]. Anyhow, that calculation
is performed more generally in the case $m_h^2\neq\mu_F^2$. This
results in extra-terms proportional to $\ln(m_h^2/\mu_F^2)$, which
vanish if, as we do, one takes the factorisation scale $\mu^2$
coinciding with the scale of the process $Q^2$. Moreover, in that case
also NLO ($\propto a_s$) appear in the matching conditions.

To summarise the PDF matching conditions at the threshold $m_h^2$, we
have that:
\begin{itemize}
\item singlet and gluon couple as follows:
\begin{equation}
{\Sigma^{(n_f+1)} \choose g^{(n_f+1)}}=\left[\begin{pmatrix} 1 & 0 \\
    0 & 1\end{pmatrix}+a_s^2A_{qq,h}^{N\!S,(2)}\begin{pmatrix} 1 & 0
    \\ 0 & 0\end{pmatrix}+a_s^2\begin{pmatrix} \tilde{A}^{S,(2)}_{hq}
    & \tilde{A}^{S,(2)}_{hg} \\A^{S,(2)}_{gq,h} &
    A_{gg,h}^{S,(2)}\end{pmatrix}\right]{\Sigma^{(n_f)} \choose
  g^{(n_f)}}\,,\label{eq:couple}
\end{equation}
\item from eqs. (\ref{valence}), (\ref{v38}), (\ref{v152435}) and (\ref{t38}), one can see that $V$, $V_{3,8,\dots,35}$ and $T_{3,8}$ behave in the same way, i.e.:
\begin{equation}
P^{(n_f+1)}=[1+a_s^2A_{qq,h}^{N\!S,(2)}]P^{(n_f)}\quad\mbox{with}\quad
P=V,V_3,\dots,V_{35},T_3,T_8\,,\label{eq:ciao}
\end{equation}
\item $T_{15}$, $T_{24}$ and $T_{35}$ have different matching
  conditions depending on the threshold. In particular: for
  $m_h^2=m_c^2$ they are given by eq. (\ref{pippo1}), for
  $m_h^2=m_b^2$ they are given by eqs. (\ref{pippo2}) and
  (\ref{pippo3}) and for $m_h^2=m_t^2$ they are given by
  eqs. (\ref{pippo4}) and (\ref{pippo5})
\end{itemize}

In the following Sections we will discuss how to write the evolution
kernels in the presence of the matching conditions. We will explicitly
consider only the forward evolution, i.e. the final scale $Q^2$
greater than the initial one $Q_0^2$. Anyway the backward evolution
($Q_0^2>Q^2$) can be easily obtained from the forward one. In fact,
given the evolution kernel $\Gamma$, the following relation holds:
\begin{equation}
\Gamma(Q^2,Q_0^2)\Gamma(Q_0^2,Q^2)=1\quad\Longrightarrow\quad\Gamma(Q^2,Q_0^2)=\Gamma^{-1}(Q_0^2,Q^2).
\end{equation}
so, if $Q^2>Q_0^2$ the code computes directly $\Gamma(Q^2,Q_0^2)$, else if $Q^2_0>Q^2$ the code evaluates first the forward evolution $\Gamma(Q_0^2,Q^2)$ and then, to get $\Gamma(Q^2,Q_0^2)$, it calculates $\Gamma^{-1}(Q_0^2,Q^2)$.

\subsection{Matching Conditions on the Evolution Kernels: 0 Thresholds Crossing}

Before to discuss the crossing of the thresholds, it would be useful
to write down the evolution kernels in the ``trivial'' situation of no
threshold crossing. There are 4 particular cases: 1) $Q_0^2<Q^2<m_c^2$
with $n_f=3$ active flavours, 2) $m_c^2<Q_0^2<Q^2<m_b^2$ with $n_f=4$
active flavours, 3) $m_b^2<Q_0^2<Q^2<m_t^2$ with $n_f=5$ active
flavours and 4) $m_t^2<Q_0^2<Q^2$ with $n_f=6$ active flavours, which
do not need the introduction of the matching conditions. So, in the
following Subsections we will write down the evolution of the whole
PDF set the these cases.

\subsubsection{$Q_0^2<Q^2<m_c^2$}
\begin{itemize}
\item \textbf{Singlet} and \textbf{gluon}:
\begin{equation}
{\Sigma^{(3)}(Q^2) \choose g^{(3)}(Q^2)} =\underbrace{\begin{pmatrix} \Gamma_{qq}& \Gamma_{qg} \\ \Gamma_{gq}& \Gamma_{gg}\end{pmatrix}}_{(Q^2,Q_0^2)}{\Sigma^{(3)}(Q_0^2) \choose g^{(3)}(Q_0^2)}
\end{equation}
\item $\mathbf{V}$:
\begin{equation}
V^{(3)}(Q^2)=\Gamma^{v}(Q^2,Q_0^2)V^{(3)}(Q^2_0)
\end{equation}
\item $\mathbf{V_3}$ and $\mathbf{V_8}$:
\begin{equation}
V^{(3)}_{3,8}(Q^2)=\Gamma^{-}(Q^2,Q_0^2)V^{(3)}_{3,8}(Q^2_0)
\end{equation}
\item $\mathbf{V_{15}}$, $\mathbf{V_{24}}$ and $\mathbf{V_{35}}$:
\begin{equation}
V_{15,24,35}^{(3)}(Q^2)=\Gamma^{v}(Q^2,Q^2_0)V^{(3)}(Q_0^2)
\end{equation}
\item $\mathbf{T_3}$ and $\mathbf{T_8}$:
\begin{equation}
T^{(3)}_{3,8}(Q^2)=\Gamma^{+}(Q^2,Q_0^2)T^{(3)}_{3,8}(Q^2_0)
\end{equation}
\item $\mathbf{T_{15}}$, $\mathbf{T_{24}}$ and $\mathbf{T_{35}}$:
\begin{equation}
T_{15,24,35}^{(3)}(Q^2) = \underbrace{\begin{pmatrix} \Gamma_{qq} & \Gamma_{qg}\end{pmatrix}}_{(Q^2,Q_0^2)}{\Sigma^{(3)}(Q_0^2) \choose g^{(3)}(Q_0^2)}
\end{equation}
\end{itemize}

\subsubsection{$m_c^2<Q_0^2<Q^2<m_b^2$}
\begin{itemize}
\item \textbf{Singlet} and \textbf{gluon}:
\begin{equation}
{\Sigma^{(4)}(Q^2) \choose g^{(4)}(Q^2)} =\underbrace{\begin{pmatrix} \Gamma_{qq}& \Gamma_{qg} \\ \Gamma_{gq}& \Gamma_{gg}\end{pmatrix}}_{(Q^2,Q_0^2)}{\Sigma^{(4)}(Q_0^2) \choose g^{(4)}(Q_0^2)}
\end{equation}
\item $\mathbf{V}$:
\begin{equation}
V^{(4)}(Q^2)=\Gamma^{v}(Q^2,Q_0^2)V^{(4)}(Q^2_0)
\end{equation}
\item $\mathbf{V_3}$, $\mathbf{V_8}$ and $\mathbf{V_{15}}$:
\begin{equation}
V^{(4)}_{3,8,15}(Q^2)=\Gamma^{-}(Q^2,Q_0^2)V^{(4)}_{3,8,15}(Q^2_0)
\end{equation}
\item $\mathbf{V_{24}}$ and $\mathbf{V_{35}}$:
\begin{equation}
V_{24,35}^{(4)}(Q^2)=\Gamma^{v}(Q^2,Q^2_0)V^{(4)}(Q_0^2)
\end{equation}
\item $\mathbf{T_3}$, $\mathbf{T_8}$ and $\mathbf{T_{15}}$:
\begin{equation}
T^{(4)}_{3,8,15}(Q^2)=\Gamma^{+}(Q^2,Q_0^2)T^{(4)}_{3,8,15}(Q^2_0)
\end{equation}
\item $\mathbf{T_{24}}$ and $\mathbf{T_{35}}$:
\begin{equation}
T_{24,35}^{(4)}(Q^2) = \underbrace{\begin{pmatrix} \Gamma_{qq} & \Gamma_{qg}\end{pmatrix}}_{(Q^2,Q_0^2)}{\Sigma^{(4)}(Q_0^2) \choose g^{(4)}(Q_0^2)}
\end{equation}
\end{itemize}

\subsubsection{$m_b^2<Q_0^2<Q^2<m_t^2$}
\begin{itemize}
\item \textbf{Singlet} and \textbf{gluon}:
\begin{equation}
{\Sigma^{(5)}(Q^2) \choose g^{(5)}(Q^2)} =\underbrace{\begin{pmatrix} \Gamma_{qq}& \Gamma_{qg} \\ \Gamma_{gq}& \Gamma_{gg}\end{pmatrix}}_{(Q^2,Q_0^2)}{\Sigma^{(5)}(Q_0^2) \choose g^{(5)}(Q_0^2)}
\end{equation}
\item $\mathbf{V}$:
\begin{equation}
V^{(5)}(Q^2)=\Gamma^{v}(Q^2,Q_0^2)V^{(5)}(Q^2_0)
\end{equation}
\item $\mathbf{V_3}$, $\mathbf{V_8}$, $\mathbf{V_{15}}$ and $\mathbf{V_{24}}$:
\begin{equation}
V^{(5)}_{3,8,15,24}(Q^2)=\Gamma^{-}(Q^2,Q_0^2)V^{(5)}_{3,8,15,24}(Q^2_0)
\end{equation}
\item $\mathbf{V_{35}}$:
\begin{equation}
V_{35}^{(5)}(Q^2)=\Gamma^{v}(Q^2,Q^2_0)V^{(5)}(Q_0^2)
\end{equation}
\item $\mathbf{T_3}$, $\mathbf{T_8}$, $\mathbf{T_{15}}$ and $\mathbf{T_{24}}$:
\begin{equation}
T^{(5)}_{3,8,15,24}(Q^2)=\Gamma^{+}(Q^2,Q_0^2)T^{(5)}_{3,8,15,24}(Q^2_0)
\end{equation}
\item $\mathbf{T_{35}}$:
\begin{equation}
T_{35}^{(5)}(Q^2) = \underbrace{\begin{pmatrix} \Gamma_{qq} & \Gamma_{qg}\end{pmatrix}}_{(Q^2,Q_0^2)}{\Sigma^{(5)}(Q_0^2) \choose g^{(5)}(Q_0^2)}
\end{equation}
\end{itemize}

\subsubsection{$m_t^2<Q_0^2<Q^2$}
\begin{itemize}
\item \textbf{Singlet} and \textbf{gluon}:
\begin{equation}
{\Sigma^{(6)}(Q^2) \choose g^{(6)}(Q^2)} =\underbrace{\begin{pmatrix} \Gamma_{qq}& \Gamma_{qg} \\ \Gamma_{gq}& \Gamma_{gg}\end{pmatrix}}_{(Q^2,Q_0^2)}{\Sigma^{(6)}(Q_0^2) \choose g^{(6)}(Q_0^2)}
\end{equation}
\item $\mathbf{V}$:
\begin{equation}
V^{(6)}(Q^2)=\Gamma^{v}(Q^2,Q_0^2)V^{(6)}(Q^2_0)
\end{equation}
\item $\mathbf{V_3}$, $\mathbf{V_8}$, $\mathbf{V_{15}}$, $\mathbf{V_{24}}$ and $\mathbf{V_{35}}$:
\begin{equation}
V^{(6)}_{3,8,15,24,35}(Q^2)=\Gamma^{-}(Q^2,Q_0^2)V^{(6)}_{3,8,15,24,35}(Q^2_0)
\end{equation}
\item $\mathbf{T_3}$, $\mathbf{T_8}$, $\mathbf{T_{15}}$, $\mathbf{T_{24}}$ and $\mathbf{T_{35}}$:
\begin{equation}
T^{(6)}_{3,8,15,24,35}(Q^2)=\Gamma^{+}(Q^2,Q_0^2)T^{(6)}_{3,8,15,24,35}(Q^2_0)
\end{equation}
\end{itemize}


\subsection{Matching Conditions on the Evolution Kernels: 1 Threshold Crossing}

In order to implement the matching conditions in our code, we will
show how to transfer them from the PDFs to the evolution kernels.

In this Section we suppose that the evolution crosses only one
threshold. We will show how the matching conditions on the PDFs modify
the form of the evolution kernels in the cases: 1)
$Q_0^2<m_c^2\leq Q^2$, 2) $Q_0^2<m_b^2\leq Q^2$ and
$Q_0^2<m_c^2\leq Q^2$.

\subsubsection{$Q_0^2<m_c^2\leq Q^2$}

In order to evolve PDFs from the scale $Q_0^2$ to $Q^2$ passing
through the threshold $m_c^2$, we have to: first evolve them from
$Q^2_0$ to $m_c^2$, where there are 3 active flavours, then increase
the number of active flavour from 3 to 4 by imposing the matching
conditions, and in the end evolve the PDFs, now having 4 active
flavours, from $m_c^2$ to the scale $Q^2$.

In what follows we will work only in the Mellin space, so we will drop
any dependence on $N$.

Let's start with singlet and gluon. We have:
\begin{equation}
{\Sigma^{(4)} \choose g^{(4)}}(Q^2) = \begin{pmatrix} \Gamma_{qq} &
  \Gamma_{qg} \\ \Gamma_{gq}&
  \Gamma_{gg}\end{pmatrix}(Q^2,m_c^2){\Sigma^{(4)} \choose
  g^{(4)}}(m_c^2)\,.
\label{couple4Qmc}
\end{equation}
From eq. (\ref{eq:couple}):
\begin{equation}
\displaystyle {\Sigma^{(4)} \choose g^{(4)}}(m_c^2)=\left[\begin{pmatrix} 1 & 0 \\ 0 & 1\end{pmatrix}+a_s^2(m_c^2)A_{qq,c}^{N\!S,(2)}\begin{pmatrix} 1 & 0 \\ 0 & 0\end{pmatrix}+a_s^2(m_c^2)\begin{pmatrix} \tilde{A}^{S,(2)}_{cq} & \tilde{A}^{S,(2)}_{cg} \\A^{S,(2)}_{gq,c} & A_{gg,c}^{S,(2)}\end{pmatrix}\right]{\Sigma^{(3)} \choose g^{(3)}}(m_c^2)
\end{equation}
now, substituting the above relation into the eq. (\ref{couple4Qmc}),
we get:
\begin{equation}
\begin{array}{l}
\displaystyle {\Sigma^{(4)} \choose g^{(4)}}(Q^2) = \begin{pmatrix} \Gamma_{qq} & \Gamma_{qg} \\ \Gamma_{gq}& \Gamma_{gg}\end{pmatrix}(Q^2,m_c^2)\Bigg[\begin{pmatrix} 1 & 0 \\ 0 & 1\end{pmatrix}+a_s^2(m_c^2)A_{qq,c}^{N\!S,(2)}\begin{pmatrix} 1 & 0 \\ 0 & 0\end{pmatrix}+\\
\\
\hspace{180pt}\displaystyle a_s^2(m_c^2)\begin{pmatrix} \tilde{A}^{S,(2)}_{cq} & \tilde{A}^{S,(2)}_{cg} \\A^{S,(2)}_{gq,c} & A_{gg,c}^{S,(2)}\end{pmatrix}\Bigg]{\Sigma^{(3)} \choose g^{(3)}}(m_c^2)\,.
\end{array}
\label{couple43Qmc}
\end{equation}
But:
\begin{equation}
{\Sigma^{(3)} \choose g^{(3)}}(m_c^2) = \begin{pmatrix} \Gamma_{qq} &
  \Gamma_{qg} \\ \Gamma_{gq}&
  \Gamma_{gg}\end{pmatrix}(m_c^2,Q_0^2){\Sigma^{(3)} \choose
  g^{(3)}}(Q_0^2)\,.
\label{couple3mcQ0}
\end{equation}
So, in the end:
\begin{equation}
\begin{array}{l}
\displaystyle {\Sigma^{(4)} \choose g^{(4)}}(Q^2) = \Bigg\{\begin{pmatrix} \Gamma_{qq} & \Gamma_{qg} \\ \Gamma_{gq}& \Gamma_{gg}\end{pmatrix}(Q^2,m_c^2)\Bigg[\begin{pmatrix} 1 & 0 \\ 0 & 1\end{pmatrix}+a_s^2(m_c^2)A_{qq,c}^{N\!S,(2)}\begin{pmatrix} 1 & 0 \\ 0 & 0\end{pmatrix}+\\
\\
\hspace{70pt}\displaystyle a_s^2(m_c^2)\begin{pmatrix} \tilde{A}^{S,(2)}_{cq} & \tilde{A}^{S,(2)}_{cg} \\A^{S,(2)}_{gq,c} & A_{gg,c}^{S,(2)}\end{pmatrix}\Bigg]\begin{pmatrix} \Gamma_{qq} & \Gamma_{qg} \\ \Gamma_{gq}& \Gamma_{gg}\end{pmatrix}(m_c^2,Q_0^2)\Bigg\}{\Sigma^{(3)} \choose g^{(3)}}(Q_0^2)
\end{array}\,.
\label{couple43QQ0}
\end{equation}

Now we consider the distributions $V$, $V_{3,8}$ and $T_{3,8}$ which
evolve re\-spe\-cti\-ve\-ly through $\Gamma^v$, $\Gamma^-$ and
$\Gamma^+$, but which obey the same matching conditions. So:
\begin{equation}
P^{(4)}(Q^2)=\Gamma^{(P)}(Q^2,m_c^2)P^{(4)}(m_c^2)
\end{equation}
so that:
\begin{equation}
P=\left\{
\begin{array}{ll}
\displaystyle V &\rightarrow \Gamma^{(P)}=\Gamma^v\\
\displaystyle V_{3,8} &\rightarrow \Gamma^{(P)}=\Gamma^-\\
\displaystyle T_{3,8} &\rightarrow \Gamma^{(P)}=\Gamma^+
\end{array}\right.
\end{equation}
but:
\begin{equation}
P^{(4)}(m_c^2)=[1+a_s^2(m_c^2)A_{qq,c}^{N\!S,(2)}]P^{(3)}(m_c^2)
\end{equation}
and:
\begin{equation}
P^{(3)}(m_c^2)=\Gamma^{(P)}(m_c^2,Q_0^2)P^{(3)}(Q^2_0)
\end{equation}
so that:
\begin{equation}
P^{(4)}(Q^2)=\left\{\Gamma^{(P)}(Q^2,m_c^2)[1+a_s^2(m_c^2)A_{qq,c}^{N\!S,(2)}]\Gamma^{(P)}(m_c^2,Q_0^2)\right\}P^{(3)}(Q^2_0)
\end{equation}

Now we consider $V_{15}$, which before $m_c^2$ evolves as:
\begin{equation}
V_{15}^{(3)}(m_c^2)=\Gamma^{v}(m_c^2,Q_0^2)V^{(3)}(Q_0^2)
\end{equation}
while after $m_c^2$ it evolves as:
\begin{equation}
V_{15}^{(4)}(Q^2)=\Gamma^{-}(Q_0^2,m_c^2)V^{(4)}_{15}(m_c^2)\,.
\end{equation}
From eq. (\ref{eq:ciao}), we find that the matching condition at $m_c^2$ is:
\begin{equation}
V^{(4)}_{15}(m_c^2)=[1+a_s^2(m_c^2)A_{qq,c}^{N\!S,(2)}] V^{(3)}_{15}(m_c^2)
\end{equation}
so:
\begin{equation}
V_{15}^{(4)}(Q^2)=\left\{\Gamma^{-}(Q^2,m_c^2)[1+a_s^2(m_c^2)A_{qq,c}^{N\!S,(2)}]\Gamma^{v}(m_c^2,Q_0^2)\right\}V^{(3)}(Q_0^2)
\end{equation}

Instead, for $Q_0^2<m_c^2\leq Q^2$, both $V_{24}$ and $V_{35}$ evolve as: 
\begin{equation}
V_{24,35}^{(4)}(Q^2)=\left\{\Gamma^{v}(Q^2,m_c^2)[1+a_s^2(m_c^2)A_{qq,c}^{N\!S,(2)}]\Gamma^{v}(m_c^2,Q_0^2)\right\}V^{(3)}(Q_0^2)
\end{equation}

Now, we consider $T_{15}$. After $m_c^2$, it evolves as:
\begin{equation}
T_{15}^{(4)}(Q^2)=\Gamma^{+}(Q^2,m_c^2)T^{(4)}_{15}(m_c^2)\,.
\end{equation}
This time the matching condition is given by the first line of eq. (\ref{pippo1}):
\begin{equation}
T_{15}^{(4)}(m_c^2)= \begin{pmatrix}  1+a_s^2(m_c^2)[A_{qq,c}^{N\!S,(2)}-3\tilde{A}^{S,(2)}_{cq}] & -3a_s^2(m_c^2)\tilde{A}^{S,(2)}_{cg} \end{pmatrix}{\Sigma^{(3)}(m_c^2) \choose g^{(3)}(m_c^2)}\,.
\end{equation}
But:
\begin{equation}
{\Sigma^{(3)}(m_c^2) \choose g^{(3)}(m_c^2) }=\begin{pmatrix} \Gamma_{qq}(m_c^2,Q_0^2) & \Gamma_{qg}(m_c^2,Q_0^2)\\ \Gamma_{gq}(m_c^2,Q_0^2) & \Gamma_{gg}(m_c^2,Q_0^2)\end{pmatrix}{\Sigma^{(3)}(Q_0^2) \choose g^{(3)}(Q_0^2)}
\end{equation}
In the end, one finds that:
\begin{equation}
\begin{array}{rcl}
\displaystyle T_{15}^{(4)}(Q^2)&=&\Bigg\{\Gamma^{+}(Q^2,m_c^2)\begin{pmatrix} 1+a_s^2(m_c^2)[A_{qq,c}^{N\!S,(2)}-3\tilde{A}^{S,(2)}_{cq}] & -3a_s^2(m_c^2)\tilde{A}^{S,(2)}_{cg}\end{pmatrix}\times\\
\\
& &\displaystyle \begin{pmatrix}\Gamma_{qq}(m_c^2,Q_0^2) & \Gamma_{qg}(m_c^2,Q_0^2)\\ \Gamma_{gq}(m_c^2,Q_0^2) & \Gamma_{gg}(m_c^2,Q_0^2)\end{pmatrix}\Bigg\}{\Sigma^{(3)}(Q_0^2) \choose g^{(3)}(Q_0^2)}
\end{array}
\end{equation} 

Now we are left only with $T_{24}$ and $T_{35}$, which evolve as the
single before and after the threshold, so they evolve exactly as the
first line of eq. (\ref{couple43QQ0}), i.e:
\begin{equation}
\begin{array}{l}
\displaystyle T_{24,35}^{(4)}(Q^2)= \Bigg\{\begin{pmatrix} \Gamma_{qq} & \Gamma_{qg} \end{pmatrix}(Q^2,m_c^2)\Bigg[\begin{pmatrix} 1 & 0 \\ 0 & 1\end{pmatrix}+a_s^2(m_c^2)A_{qq,c}^{N\!S,(2)}\begin{pmatrix} 1 & 0 \\ 0 & 0\end{pmatrix}+\\
\\
\hspace{70pt}\displaystyle a_s^2(m_c^2)\begin{pmatrix} \tilde{A}^{S,(2)}_{cq} & \tilde{A}^{S,(2)}_{cg} \\A^{S,(2)}_{gq,c} & A_{gg,c}^{S,(2)}\end{pmatrix}\Bigg]\begin{pmatrix} \Gamma_{qq} & \Gamma_{qg} \\ \Gamma_{gq}& \Gamma_{gg}\end{pmatrix}(m_c^2,Q_0^2)\Bigg\}{\Sigma^{(3)} \choose g^{(3)}}(Q_0^2)
\end{array}\,.
\end{equation}

Now, let us summarise what happens to the evolution kernels, by
introducing the matching conditions, if one crosses the $m_c^2$
threshold. We remind that the matching conditions appear only from the
NNLO.
\begin{itemize}
\item \textbf{Singlet} and \textbf{gluon}:
\begin{equation}
{\Sigma^{(4)}(Q^2) \choose g^{(4)}(Q^2)} = \Bigg\{\underbrace{\begin{pmatrix} \Gamma_{qq} & \Gamma_{qg} \\ \Gamma_{gq}& \Gamma_{gg}\end{pmatrix}}_{(Q^2,m_c^2)}\begin{pmatrix} M_{11}^c & M_{12}^c \\ M_{21}^c & M_{22}^c\end{pmatrix}\underbrace{\begin{pmatrix} \Gamma_{qq}& \Gamma_{qg} \\ \Gamma_{gq}& \Gamma_{gg}\end{pmatrix}}_{(m_c^2,Q_0^2)}\Bigg\}{\Sigma^{(3)}(Q_0^2) \choose g^{(3)}(Q_0^2)}
\end{equation}
where:
\begin{equation}
\begin{pmatrix} M_{11}^c & M_{12}^c \\ M_{21}^c & M_{22}^c\end{pmatrix}=\begin{pmatrix} 1 & 0 \\ 0 & 1\end{pmatrix}+a_s^2(m_c^2)A_{qq,c}^{N\!S,(2)}\begin{pmatrix} 1 & 0 \\ 0 & 0\end{pmatrix}+a_s^2(m_c^2)\begin{pmatrix} \tilde{A}^{S,(2)}_{cq} & \tilde{A}^{S,(2)}_{cg} \\A^{S,(2)}_{gq,c} & A_{gg,c}^{S,(2)}\end{pmatrix}
\end{equation}
\item $\mathbf{V}$:
\begin{equation}
V^{(4)}(Q^2)=\left\{\Gamma^{v}(Q^2,m_c^2)[1+a_s^2(m_c^2)A_{qq,c}^{N\!S,(2)}]\Gamma^{v}(m_c^2,Q_0^2)\right\}V^{(3)}(Q^2_0)
\end{equation}
\item $\mathbf{V_3}$ and $\mathbf{V_8}$:
\begin{equation}
V^{(4)}_{3,8}(Q^2)=\left\{\Gamma^{-}(Q^2,m_c^2)[1+a_s^2(m_c^2)A_{qq,c}^{N\!S,(2)}]\Gamma^{-}(m_c^2,Q_0^2)\right\}V^{(3)}_{3,8}(Q^2_0)
\end{equation}
\item $\mathbf{V_{15}}$:
\begin{equation}
V_{15}^{(4)}(Q^2)=\left\{\Gamma^{-}(Q^2,m_c^2)[1+a_s^2(m_c^2)A_{qq,c}^{N\!S,(2)}]\Gamma^{v}(m_c^2,Q_0^2)\right\}V^{(3)}(Q_0^2)
\end{equation}
\item $\mathbf{V_{24}}$ and $\mathbf{V_{35}}$:
\begin{equation}
V_{24,35}^{(4)}(Q^2)=\left\{\Gamma^{v}(Q^2,m_c^2)[1+a_s^2(m_c^2)A_{qq,c}^{N\!S,(2)}]\Gamma^{v}(m_c^2,Q_0^2)\right\}V^{(3)}(Q_0^2)
\end{equation}
\item $\mathbf{T_3}$ and $\mathbf{T_8}$:
\begin{equation}
T^{(4)}_{3,8}(Q^2)=\left\{\Gamma^{+}(Q^2,m_c^2)[1+a_s^2(m_c^2)A_{qq,c}^{N\!S,(2)}]\Gamma^{+}(m_c^2,Q_0^2)\right\}T^{(3)}_{3,8}(Q^2_0)
\end{equation}
\item $\mathbf{T_{15}}$:
\begin{equation}
\begin{array}{rcl}
\displaystyle T_{15}^{(4)}(Q^2)&=&\Bigg\{\Gamma^{+}(Q^2,m_c^2)\begin{pmatrix} 1+a_s^2(m_c^2)[A_{qq,c}^{N\!S,(2)}-3\tilde{A}^{S,(2)}_{cq}] & -3a_s^2(m_c^2)\tilde{A}^{S,(2)}_{cg}\end{pmatrix}\times\\
\\
 & & \displaystyle \underbrace{\begin{pmatrix}\Gamma_{qq}& \Gamma_{qg} \\ \Gamma_{gq} & \Gamma_{gg}\end{pmatrix}}_{(m_c^2,Q_0^2)}\Bigg\}{\Sigma^{(3)}(Q_0^2) \choose g^{(3)}(Q_0^2)}
\end{array}
\end{equation} 
\item $\mathbf{T_{24}}$ and $\mathbf{T_{35}}$:
\begin{equation}
T_{24,35}^{(4)}(Q^2) = \Bigg\{\underbrace{\begin{pmatrix} \Gamma_{qq} & \Gamma_{qg}\end{pmatrix}}_{(Q^2,m_c^2)}\begin{pmatrix} M_{11}^c & M_{12}^c \\ M_{21}^c & M_{22}^c\end{pmatrix}\underbrace{\begin{pmatrix} \Gamma_{qq}& \Gamma_{qg} \\ \Gamma_{gq}& \Gamma_{gg}\end{pmatrix}}_{(m_c^2,Q_0^2)}\Bigg\}{\Sigma^{(3)}(Q_0^2) \choose g^{(3)}(Q_0^2)}
\end{equation}
\end{itemize}

Now, it is very easy to rewrite the above summary for the crossing of
the remaining thresholds $m_b^2$ ($Q_0^2<m_b^2\leq Q^2$) and $m_t^2$
($Q_0^2<m_t^2\leq Q^2$).

\subsubsection{$Q_0^2<m_b^2\leq Q^2$}
\begin{itemize}
\item \textbf{Singlet} and \textbf{gluon}:
\begin{equation}
{\Sigma^{(5)}(Q^2) \choose g^{(5)}(Q^2)} = \Bigg\{\underbrace{\begin{pmatrix} \Gamma_{qq} & \Gamma_{qg} \\ \Gamma_{gq}& \Gamma_{gg}\end{pmatrix}}_{(Q^2,m_b^2)}\begin{pmatrix} M_{11}^b & M_{12}^b \\ M_{21}^b & M_{22}^b\end{pmatrix}\underbrace{\begin{pmatrix} \Gamma_{qq}& \Gamma_{qg} \\ \Gamma_{gq}& \Gamma_{gg}\end{pmatrix}}_{(m_b^2,Q_0^2)}\Bigg\}{\Sigma^{(4)}(Q_0^2) \choose g^{(4)}(Q_0^2)}
\end{equation}
where:
\begin{equation}
\begin{pmatrix} M_{11}^b & M_{12}^b \\ M_{21}^b & M_{22}^b\end{pmatrix}=\begin{pmatrix} 1 & 0 \\ 0 & 1\end{pmatrix}+a_s^2(m_b^2)A_{qq,b}^{N\!S,(2)}\begin{pmatrix} 1 & 0 \\ 0 & 0\end{pmatrix}+a_s^2(m_b^2)\begin{pmatrix} \tilde{A}^{S,(2)}_{bq} & \tilde{A}^{S,(2)}_{bg} \\A^{S,(2)}_{gq,b} & A_{gg,b}^{S,(2)}\end{pmatrix}
\end{equation}
\item $\mathbf{V}$:
\begin{equation}
V^{(5)}(Q^2)=\left\{\Gamma^{v}(Q^2,m_b^2)[1+a_s^2(m_b^2)A_{qq,b}^{N\!S,(2)}]\Gamma^{v}(m_b^2,Q_0^2)\right\}V^{(4)}(Q^2_0)
\end{equation}
\item $\mathbf{V_3}$, $\mathbf{V_8}$ and $\mathbf{V_{15}}$:
\begin{equation}
V^{(5)}_{3,8,15}(Q^2)=\left\{\Gamma^{-}(Q^2,m_b^2)[1+a_s^2(m_b^2)A_{qq,b}^{N\!S,(2)}]\Gamma^{-}(m_b^2,Q_0^2)\right\}V^{(4)}_{3,8,15}(Q^2_0)
\end{equation}
\item $\mathbf{V_{24}}$:
\begin{equation}
V_{24}^{(5)}(Q^2)=\left\{\Gamma^{-}(Q_0^2,m_b^2)[1+a_s^2(m_b^2)A_{qq,b}^{N\!S,(2)}]\Gamma^{v}(m_b^2,Q_0^2)\right\}V^{(4)}(Q_0^2)
\end{equation}
\item $\mathbf{V_{35}}$:
\begin{equation}
V_{35}^{(5)}(Q^2)=\left\{\Gamma^{v}(Q_0^2,m_b^2)[1+a_s^2(m_b^2)A_{qq,b}^{N\!S,(2)}]\Gamma^{v}(m_b^2,Q_0^2)\right\}V^{(4)}(Q_0^2)
\end{equation}
\item $\mathbf{T_3}$, $\mathbf{T_8}$ and $\mathbf{T_{15}}$:
\begin{equation}
T^{(5)}_{3,8,15}(Q^2)=\left\{\Gamma^{+}(Q^2,m_b^2)[1+a_s^2(m_b^2)A_{qq,b}^{N\!S,(2)}]\Gamma^{+}(m_b^2,Q_0^2)\right\}T^{(4)}_{3,8,15}(Q^2_0)
\end{equation}
\item $\mathbf{T_{24}}$:
\begin{equation}
\begin{array}{rcl}
\displaystyle T_{24}^{(5)}(Q^2)&=&\Bigg\{\Gamma^{+}(Q^2,m_b^2)\begin{pmatrix} 1+a_s^2(m_b^2)[A_{qq,b}^{N\!S,(2)}-4\tilde{A}^{S,(2)}_{bq}] & -4a_s^2(m_b^2)\tilde{A}^{S,(2)}_{bg}\end{pmatrix}\times\\
\\
 & & \displaystyle \underbrace{\begin{pmatrix}\Gamma_{qq}& \Gamma_{qg} \\ \Gamma_{gq} & \Gamma_{gg}\end{pmatrix}}_{(m_b^2,Q_0^2)}\Bigg\}{\Sigma^{(4)}(Q_0^2) \choose g^{(4)}(Q_0^2)}
\end{array}
\end{equation} 
\item $\mathbf{T_{35}}$:
\begin{equation}
T_{35}^{(5)}(Q^2) = \Bigg\{\underbrace{\begin{pmatrix} \Gamma_{qq} & \Gamma_{qg}\end{pmatrix}}_{(Q^2,m_b^2)}\begin{pmatrix} M_{11}^b & M_{12}^b \\ M_{21}^b & M_{22}^b\end{pmatrix}\underbrace{\begin{pmatrix} \Gamma_{qq}& \Gamma_{qg} \\ \Gamma_{gq}& \Gamma_{gg}\end{pmatrix}}_{(m_b^2,Q_0^2)}\Bigg\}{\Sigma^{(4)}(Q_0^2) \choose g^{(4)}(Q_0^2)}
\end{equation}
\end{itemize}

\subsubsection{$Q_0^2<m_t^2\leq Q^2$}
\begin{itemize}
\item \textbf{Singlet} and \textbf{gluon}:
\begin{equation}
{\Sigma^{(6)}(Q^2) \choose g^{(6)}(Q^2)} = \Bigg\{\underbrace{\begin{pmatrix} \Gamma_{qq} & \Gamma_{qg} \\ \Gamma_{gq}& \Gamma_{gg}\end{pmatrix}}_{(Q^2,m_t^2)}\begin{pmatrix} M_{11}^t & M_{12}^t \\ M_{21}^t & M_{22}^t\end{pmatrix}\underbrace{\begin{pmatrix} \Gamma_{qq}& \Gamma_{qg} \\ \Gamma_{gq}& \Gamma_{gg}\end{pmatrix}}_{(m_t^2,Q_0^2)}\Bigg\}{\Sigma^{(5)}(Q_0^2) \choose g^{(5)}(Q_0^2)}
\end{equation}
where:
\begin{equation}
\begin{pmatrix} M_{11}^t & M_{12}^t \\ M_{21}^t & M_{22}^t\end{pmatrix}=\begin{pmatrix} 1 & 0 \\ 0 & 1\end{pmatrix}+a_s^2(m_t^2)A_{qq,t}^{N\!S,(2)}\begin{pmatrix} 1 & 0 \\ 0 & 0\end{pmatrix}+a_s^2(m_t^2)\begin{pmatrix} \tilde{A}^{S,(2)}_{tq} & \tilde{A}^{S,(2)}_{tg} \\A^{S,(2)}_{gq,t} & A_{gg,t}^{S,(2)}\end{pmatrix}
\end{equation}
\item $\mathbf{V}$:
\begin{equation}
V^{(6)}(Q^2)=\left\{\Gamma^{v}(Q^2,m_t^2)[1+a_s^2(m_t^2)A_{qq,t}^{N\!S,(2)}]\Gamma^{v}(m_t^2,Q_0^2)\right\}V^{(5)}(Q^2_0)
\end{equation}
\item $\mathbf{V_3}$, $\mathbf{V_8}$, $\mathbf{V_{15}}$ and $\mathbf{V_{24}}$:
\begin{equation}
V^{(6)}_{3,8,15,24}(Q^2)=\left\{\Gamma^{-}(Q^2,m_t^2)[1+a_s^2(m_t^2)A_{qq,t}^{N\!S,(2)}]\Gamma^{-}(m_t^2,Q_0^2)\right\}V^{(5)}_{3,8,15,24}(Q^2_0)
\end{equation}
\item $\mathbf{V_{35}}$:
\begin{equation}
V_{35}^{(6)}(Q^2)=\left\{\Gamma^{-}(Q_0^2,m_t^2)[1+a_s^2(m_t^2)A_{qq,t}^{N\!S,(2)}]\Gamma^{v}(m_t^2,Q_0^2)\right\}V^{(5)}(Q_0^2)
\end{equation}
\item $\mathbf{T_3}$, $\mathbf{T_8}$, $\mathbf{T_{15}}$ and $\mathbf{T_{24}}$:
\begin{equation}
T^{(6)}_{3,8,15,24}(Q^2)=\left\{\Gamma^{+}(Q^2,m_t^2)[1+a_s^2(m_t^2)A_{qq,t}^{N\!S,(2)}]\Gamma^{+}(m_t^2,Q_0^2)\right\}T^{(5)}_{3,8,15,24}(Q^2_0)
\end{equation}
\item $\mathbf{T_{35}}$:
\begin{equation}
\begin{array}{rcl}
\displaystyle T_{35}^{(6)}(Q^2)&=&\Bigg\{\Gamma^{+}(Q^2,m_t^2)\begin{pmatrix} 1+a_s^2(m_t^2)[A_{qq,t}^{N\!S,(2)}-5\tilde{A}^{S,(2)}_{tq}] & -5a_s^2(m_t^2)\tilde{A}^{S,(2)}_{tg}\end{pmatrix}\times\\
\\
 & & \displaystyle \underbrace{\begin{pmatrix}\Gamma_{qq}& \Gamma_{qg} \\ \Gamma_{gq} & \Gamma_{gg}\end{pmatrix}}_{(m_t^2,Q_0^2)}\Bigg\}{\Sigma^{(5)}(Q_0^2) \choose g^{(5)}(Q_0^2)}
\end{array}
\end{equation} 
\end{itemize}

\subsection{Matching Conditions on the Evolution Kernels: 2 Thresholds Crossing}

In this Section we will discuss the case in which the evolution
crosses two thresholds. Therefore there are only two situations: 1)
$Q_0^2<m_c^2<m_b^2\leq Q^2$ and 2) $Q_0^2<m_b^2<m_t^2\leq Q^2$.
Anyway, there is nothing new, indeed to obtain such evolution kernels
we have just to ``merge'' together what we have already done in the
previous Section.

\subsubsection{$Q_0^2<m_c^2<m_b^2\leq Q^2$}

\begin{itemize}
\item \textbf{Singlet} and \textbf{gluon}:
\begin{equation}
\begin{array}{rcl}
\displaystyle {\Sigma^{(5)}(Q^2) \choose g^{(5)}(Q^2)} &=& \displaystyle \Bigg\{\underbrace{\begin{pmatrix} \Gamma_{qq} & \Gamma_{qg} \\ \Gamma_{gq}& \Gamma_{gg}\end{pmatrix}}_{(Q^2,m_b^2)}\begin{pmatrix} M_{11}^b & M_{12}^b \\ M_{21}^b & M_{22}^b\end{pmatrix}\underbrace{\begin{pmatrix} \Gamma_{qq}& \Gamma_{qg} \\ \Gamma_{gq}& \Gamma_{gg}\end{pmatrix}}_{(m_b^2,m_c^2)}\times\\
\\
 & &\displaystyle\begin{pmatrix} M_{11}^c & M_{12}^c \\ M_{21}^c & M_{22}^c\end{pmatrix}\underbrace{\begin{pmatrix} \Gamma_{qq}& \Gamma_{qg} \\ \Gamma_{gq}& \Gamma_{gg}\end{pmatrix}}_{(m_c^2,Q_0^2)}\Bigg\}{\Sigma^{(3)}(Q_0^2) \choose g^{(3)}(Q_0^2)}
\end{array}
\end{equation}
\item $\mathbf{V}$:
\begin{equation}
\begin{array}{rcl}
V^{(5)}(Q^2)&=&\displaystyle \Big\{\Gamma^{v}(Q^2,m_b^2)[1+a_s^2(m_b^2)A_{qq,b}^{N\!S,(2)}]\times\\
\\
 & &\displaystyle \Gamma^{v}(m_b^2,m_c^2)[1+a_s^2(m_c^2)A_{qq,c}^{N\!S,(2)}]\Gamma^{v}(m_c^2,Q_0^2)\Big\}V^{(3)}(Q^2_0)
\end{array}
\end{equation}
\item $\mathbf{V_3}$ and $\mathbf{V_8}$:
\begin{equation}
\begin{array}{rcl}
V^{(5)}_{3,8}(Q^2)&=&\displaystyle \Big\{\Gamma^{-}(Q^2,m_b^2)[1+a_s^2(m_b^2)A_{qq,b}^{N\!S,(2)}]\times\\
\\
& & \displaystyle \Gamma^{-}(m_b^2,m_c^2)[1+a_s^2(m_c^2)A_{qq,c}^{N\!S,(2)}]\Gamma^{-}(m_c^2,Q_0^2)\Big\}V^{(3)}_{3,8}(Q^2_0)
\end{array}
\end{equation}
\item $\mathbf{V_{15}}$:
\begin{equation}
\begin{array}{rcl}
V^{(5)}_{15}(Q^2)&=&\displaystyle \Big\{\Gamma^{-}(Q^2,m_b^2)[1+a_s^2(m_b^2)A_{qq,b}^{N\!S,(2)}]\times\\
\\
& & \displaystyle \Gamma^{-}(m_b^2,m_c^2)[1+a_s^2(m_c^2)A_{qq,c}^{N\!S,(2)}]\Gamma^{v}(m_c^2,Q_0^2)\Big\}V^{(3)}(Q^2_0)
\end{array}
\end{equation}
\item $\mathbf{V_{24}}$:
\begin{equation}
\begin{array}{rcl}
V^{(5)}_{24}(Q^2)&=&\displaystyle \Big\{\Gamma^{-}(Q^2,m_b^2)[1+a_s^2(m_b^2)A_{qq,b}^{N\!S,(2)}]\times\\
\\
& & \displaystyle \Gamma^{v}(m_b^2,m_c^2)[1+a_s^2(m_c^2)A_{qq,c}^{N\!S,(2)}]\Gamma^{v}(m_c^2,Q_0^2)\Big\}V^{(3)}(Q^2_0)
\end{array}
\end{equation}
\item $\mathbf{V_{35}}$:
\begin{equation}
\begin{array}{rcl}
V^{(5)}_{35}(Q^2)&=&\displaystyle \Big\{\Gamma^{v}(Q^2,m_b^2)[1+a_s^2(m_b^2)A_{qq,b}^{N\!S,(2)}]\times\\
\\
& & \displaystyle \Gamma^{v}(m_b^2,m_c^2)[1+a_s^2(m_c^2)A_{qq,c}^{N\!S,(2)}]\Gamma^{v}(m_c^2,Q_0^2)\Big\}V^{(3)}(Q^2_0)
\end{array}
\end{equation}

\item $\mathbf{T_3}$ and $\mathbf{T_8}$:
\begin{equation}
\begin{array}{rcl}
T^{(5)}_{3,8}(Q^2)&=&\displaystyle \Big\{\Gamma^{+}(Q^2,m_b^2)[1+a_s^2(m_b^2)A_{qq,b}^{N\!S,(2)}]\times\\
\\
& & \displaystyle \Gamma^{+}(m_b^2,m_c^2)[1+a_s^2(m_c^2)A_{qq,c}^{N\!S,(2)}]\Gamma^{+}(m_c^2,Q_0^2)\Big\}T^{(3)}_{3,8}(Q^2_0)
\end{array}
\end{equation}

\item $\mathbf{T_{15}}$:
\begin{equation}
\begin{array}{rcl}
T^{(5)}_{15}(Q^2)&=&\displaystyle \Bigg\{\Gamma^{+}(Q^2,m_b^2)[1+a_s^2(m_b^2)A_{qq,b}^{N\!S,(2)}]\Gamma^{+}(m_b^2,m_c^2)\times\\
\\
& & \displaystyle \begin{pmatrix} 1+a_s^2(m_c^2)[A_{qq,c}^{N\!S,(2)}-3\tilde{A}^{S,(2)}_{cq}] & -3a_s^2(m_c^2)\tilde{A}^{S,(2)}_{cg}\end{pmatrix}\underbrace{\begin{pmatrix}\Gamma_{qq}& \Gamma_{qg} \\ \Gamma_{gq} & \Gamma_{gg}\end{pmatrix}}_{(m_c^2,Q_0^2)}\Bigg\}{\Sigma^{(3)}(Q_0^2) \choose g^{(3)}(Q_0^2)}
\end{array}
\end{equation}

\item $\mathbf{T_{24}}$:
\begin{equation}
\begin{array}{rcl}
\displaystyle T_{24}^{(5)}(Q^2)&=&\Bigg\{\Gamma^{+}(Q^2,m_b^2)\begin{pmatrix} 1+a_s^2(m_b^2)[A_{qq,b}^{N\!S,(2)}-4\tilde{A}^{S,(2)}_{bq}] & -4a_s^2(m_b^2)\tilde{A}^{S,(2)}_{bg}\end{pmatrix}\times\\
\\
 & & \displaystyle \underbrace{\begin{pmatrix}\Gamma_{qq}& \Gamma_{qg} \\ \Gamma_{gq} & \Gamma_{gg}\end{pmatrix}}_{(m_b^2,m_c^2)}\begin{pmatrix} M^c_{11} & M^c_{12} \\ M_{21}^c & M_{22}^c \end{pmatrix}\underbrace{\begin{pmatrix}\Gamma_{qq}& \Gamma_{qg} \\ \Gamma_{gq} & \Gamma_{gg}\end{pmatrix}}_{(m_c^2,Q_0^2)}\Bigg\}{\Sigma^{(3)}(Q_0^2) \choose g^{(3)}(Q_0^2)}
\end{array}
\end{equation}
\item $\mathbf{T_{35}}$:
\begin{equation}
\begin{array}{rcl}
T_{35}^{(5)}(Q^2) &=&\displaystyle \Bigg\{\underbrace{\begin{pmatrix} \Gamma_{qq} & \Gamma_{qg}\end{pmatrix}}_{(Q^2,m_b^2)}\begin{pmatrix} M_{11}^b & M_{12}^b \\ M_{21}^b & M_{22}^b\end{pmatrix}\underbrace{\begin{pmatrix} \Gamma_{qq} & \Gamma_{qg} \\ \Gamma_{gq}& \Gamma_{gg}\end{pmatrix}}_{(m_b^2,m_c^2)}\times\\
\\
& & \displaystyle \begin{pmatrix} M^c_{11} & M^c_{12} \\ M_{21}^c & M_{22}^c \end{pmatrix}\underbrace{\begin{pmatrix}\Gamma_{qq}& \Gamma_{qg} \\ \Gamma_{gq} & \Gamma_{gg}\end{pmatrix}}_{(m_c^2,Q_0^2)}\Bigg\} {\Sigma^{(3)}(Q_0^2) \choose g^{(3)}(Q_0^2)}
\end{array}
\end{equation}
\end{itemize}

\subsubsection{$Q_0^2<m_b^2<m_t^2\leq Q^2$}

\begin{itemize}
\item \textbf{Singlet} and \textbf{gluon}:
\begin{equation}
\begin{array}{rcl}
\displaystyle {\Sigma^{(6)}(Q^2) \choose g^{(6)}(Q^2)} &=& \displaystyle \Bigg\{\underbrace{\begin{pmatrix} \Gamma_{qq} & \Gamma_{qg} \\ \Gamma_{gq}& \Gamma_{gg}\end{pmatrix}}_{(Q^2,m_t^2)}\begin{pmatrix} M_{11}^t & M_{12}^t \\ M_{21}^t & M_{22}^t\end{pmatrix}\underbrace{\begin{pmatrix} \Gamma_{qq}& \Gamma_{qg} \\ \Gamma_{gq}& \Gamma_{gg}\end{pmatrix}}_{(m_t^2,m_b^2)}\times\\
\\
 & &\displaystyle\begin{pmatrix} M_{11}^b & M_{12}^b \\ M_{21}^b & M_{22}^b\end{pmatrix}\underbrace{\begin{pmatrix} \Gamma_{qq}& \Gamma_{qg} \\ \Gamma_{gq}& \Gamma_{gg}\end{pmatrix}}_{(m_b^2,Q_0^2)}\Bigg\}{\Sigma^{(4)}(Q_0^2) \choose g^{(4)}(Q_0^2)}
\end{array}
\end{equation}
\item $\mathbf{V}$:
\begin{equation}
\begin{array}{rcl}
V^{(6)}(Q^2)&=&\displaystyle \Big\{\Gamma^{v}(Q^2,m_t^2)[1+a_s^2(m_t^2)A_{qq,t}^{N\!S,(2)}]\times\\
\\
 & &\displaystyle \Gamma^{v}(m_t^2,m_b^2)[1+a_s^2(m_b^2)A_{qq,b}^{N\!S,(2)}]\Gamma^{v}(m_b^2,Q_0^2)\Big\}V^{(4)}(Q^2_0)
\end{array}
\end{equation}
\item $\mathbf{V_3}$, $\mathbf{V_8}$ and $\mathbf{V_{15}}$:
\begin{equation}
\begin{array}{rcl}
V^{(6)}_{3,8,15}(Q^2)&=&\displaystyle \Big\{\Gamma^{-}(Q^2,m_t^2)[1+a_s^2(m_t^2)A_{qq,t}^{N\!S,(2)}]\times\\
\\
& & \displaystyle \Gamma^{-}(m_t^2,m_b^2)[1+a_s^2(m_b^2)A_{qq,b}^{N\!S,(2)}]\Gamma^{-}(m_b^2,Q_0^2)\Big\}V^{(4)}_{3,8,15}(Q^2_0)
\end{array}
\end{equation}
\item $\mathbf{V_{24}}$:
\begin{equation}
\begin{array}{rcl}
V^{(6)}_{24}(Q^2)&=&\displaystyle \Big\{\Gamma^{-}(Q^2,m_t^2)[1+a_s^2(m_t^2)A_{qq,t}^{N\!S,(2)}]\times\\
\\
& & \displaystyle \Gamma^{-}(m_t^2,m_b^2)[1+a_s^2(m_b^2)A_{qq,b}^{N\!S,(2)}]\Gamma^{v}(m_b^2,Q_0^2)\Big\}V^{(4)}(Q^2_0)
\end{array}
\end{equation}
\item $\mathbf{V_{35}}$:
\begin{equation}
\begin{array}{rcl}
V^{(6)}_{35}(Q^2)&=&\displaystyle \Big\{\Gamma^{-}(Q^2,m_t^2)[1+a_s^2(m_t^2)A_{qq,t}^{N\!S,(2)}]\times\\
\\
& & \displaystyle \Gamma^{v}(m_t^2,m_b^2)[1+a_s^2(m_b^2)A_{qq,b}^{N\!S,(2)}]\Gamma^{v}(m_b^2,Q_0^2)\Big\}V^{(4)}(Q^2_0)
\end{array}
\end{equation}

\item $\mathbf{T_3}$, $\mathbf{T_8}$ and $\mathbf{T_{15}}$:
\begin{equation}
\begin{array}{rcl}
T^{(6)}_{3,8,15}(Q^2)&=&\displaystyle \Big\{\Gamma^{+}(Q^2,m_t^2)[1+a_s^2(m_t^2)A_{qq,t}^{N\!S,(2)}]\times\\
\\
& & \displaystyle \Gamma^{+}(m_t^2,m_b^2)[1+a_s^2(m_b^2)A_{qq,b}^{N\!S,(2)}]\Gamma^{+}(m_b^2,Q_0^2)\Big\}T^{(4)}_{3,8,15}(Q^2_0)
\end{array}
\end{equation}

\item $\mathbf{T_{24}}$:
\begin{equation}
\begin{array}{rcl}
T^{(6)}_{24}(Q^2)&=&\displaystyle \Bigg\{\Gamma^{+}(Q^2,m_t^2)[1+a_s^2(m_t^2)A_{qq,t}^{N\!S,(2)}]\Gamma^{+}(m_t^2,m_b^2)\times\\
\\
& & \displaystyle \begin{pmatrix} 1+a_s^2(m_b^2)[A_{qq,b}^{N\!S,(2)}-4\tilde{A}^{S,(2)}_{bq}] & -4a_s^2(m_b^2)\tilde{A}^{S,(2)}_{bg}\end{pmatrix}\underbrace{\begin{pmatrix}\Gamma_{qq}& \Gamma_{qg} \\ \Gamma_{gq} & \Gamma_{gg}\end{pmatrix}}_{(m_b^2,Q_0^2)}\Bigg\}{\Sigma^{(4)}(Q_0^2) \choose g^{(4)}(Q_0^2)}
\end{array}
\end{equation}

\item $\mathbf{T_{35}}$:
\begin{equation}
\begin{array}{rcl}
\displaystyle T_{35}^{(6)}(Q^2)&=&\Bigg\{\Gamma^{+}(Q^2,m_t^2)\begin{pmatrix} 1+a_s^2(m_t^2)[A_{qq,t}^{N\!S,(2)}-5\tilde{A}^{S,(2)}_{tq}] & -5a_s^2(m_t^2)\tilde{A}^{S,(2)}_{tg}\end{pmatrix}\times\\
\\
 & & \displaystyle \underbrace{\begin{pmatrix}\Gamma_{qq}& \Gamma_{qg} \\ \Gamma_{gq} & \Gamma_{gg}\end{pmatrix}}_{(m_t^2,m_b^2)}\begin{pmatrix} M^b_{11} & M^b_{12} \\ M_{21}^b & M_{22}^b \end{pmatrix}\underbrace{\begin{pmatrix}\Gamma_{qq}& \Gamma_{qg} \\ \Gamma_{gq} & \Gamma_{gg}\end{pmatrix}}_{(m_b^2,Q_0^2)}\Bigg\}{\Sigma^{(4)}(Q_0^2) \choose g^{(4)}(Q_0^2)}
\end{array}
\end{equation}
\end{itemize}

\subsection{Matching Conditions on the Evolution Kernels: 3 Thresholds Crossing}

In this Section we will discuss the case in which the evolution
crosses three thresholds. Therefore there ais only one situation:
$Q_0^2<m_c^2<m_b^2,m_t^2\leq Q^2$

\subsubsection{$Q_0^2<m_c^2<m_b^2<m_t^2\leq Q^2$}
\begin{itemize}
\item \textbf{Singlet} and \textbf{gluon}:
\begin{equation}
\begin{array}{rcl}
\displaystyle {\Sigma^{(6)}(Q^2) \choose g^{(6)}(Q^2)} &=& \displaystyle \Bigg\{\underbrace{\begin{pmatrix} \Gamma_{qq} & \Gamma_{qg} \\ \Gamma_{gq}& \Gamma_{gg}\end{pmatrix}}_{(Q^2,m_t^2)}\begin{pmatrix} M_{11}^t & M_{12}^t \\ M_{21}^t & M_{22}^t\end{pmatrix}\underbrace{\begin{pmatrix} \Gamma_{qq}& \Gamma_{qg} \\ \Gamma_{gq}& \Gamma_{gg}\end{pmatrix}}_{(m_t^2,m_b^2)}\times\\
\\
 & &\displaystyle\begin{pmatrix} M_{11}^b & M_{12}^b \\ M_{21}^b & M_{22}^b\end{pmatrix}\underbrace{\begin{pmatrix} \Gamma_{qq}& \Gamma_{qg} \\ \Gamma_{gq}& \Gamma_{gg}\end{pmatrix}}_{(m_b^2,m_c^2)}\begin{pmatrix} M_{11}^c & M_{12}^c \\ M_{21}^c & M_{22}^c\end{pmatrix}\\
\\
& &\displaystyle \underbrace{\begin{pmatrix} \Gamma_{qq}& \Gamma_{qg} \\ \Gamma_{gq}& \Gamma_{gg}\end{pmatrix}}_{(m_c^2,Q_0^2)}\Bigg\}{\Sigma^{(3)}(Q_0^2) \choose g^{(3)}(Q_0^2)}
\end{array}
\end{equation}
\item $\mathbf{V}$:
\begin{equation}
\begin{array}{rcl}
V^{(6)}(Q^2)&=&\displaystyle \Big\{\Gamma^{v}(Q^2,m_t^2)[1+a_s^2(m_t^2)A_{qq,t}^{N\!S,(2)}]\times\\
\\
 & &\displaystyle \Gamma^{v}(m_t^2,m_b^2)[1+a_s^2(m_b^2)A_{qq,b}^{N\!S,(2)}]\Gamma^{v}(m_b^2,m_c^2)\\
\\
& & \displaystyle[1+a_s^2(m_c^2)A_{qq,c}^{N\!S,(2)}]\Gamma^{v}(m_c^2,Q_0^2)\Big\}V^{(3)}(Q^2_0)
\end{array}
\end{equation}
\item $\mathbf{V_3}$ and $\mathbf{V_8}$:
\begin{equation}
\begin{array}{rcl}
V^{(6)}_{3,8}(Q^2)&=&\displaystyle \Big\{\Gamma^{-}(Q^2,m_t^2)[1+a_s^2(m_t^2)A_{qq,t}^{N\!S,(2)}]\times\\
\\
& & \displaystyle \Gamma^{-}(m_t^2,m_b^2)[1+a_s^2(m_b^2)A_{qq,b}^{N\!S,(2)}]\Gamma^{-}(m_b^2,m_c^2)\\
\\
& & \displaystyle [1+a_s^2(m_c^2)A_{qq,c}^{N\!S,(2)}]\Gamma^{-}(m_c^2,Q_0^2)\Big\}V^{(3)}_{3,8}(Q^2_0)
\end{array}
\end{equation}
\item $\mathbf{V_{15}}$:
\begin{equation}
\begin{array}{rcl}
V^{(6)}_{15}(Q^2)&=&\displaystyle \Big\{\Gamma^{-}(Q^2,m_t^2)[1+a_s^2(m_t^2)A_{qq,t}^{N\!S,(2)}]\times\\
\\
& & \displaystyle \Gamma^{-}(m_t^2,m_b^2)[1+a_s^2(m_b^2)A_{qq,b}^{N\!S,(2)}]\Gamma^{-}(m_b^2,m_c^2)\\ 
\\
& & \displaystyle [1+a_s^2(m_c^2)A_{qq,c}^{N\!S,(2)}]\Gamma^{v}(m_c^2,Q_0^2)\Big\}V^{(3)}(Q^2_0)
\end{array}
\end{equation}
\item $\mathbf{V_{24}}$:
\begin{equation}
\begin{array}{rcl}
V^{(6)}_{15}(Q^2)&=&\displaystyle \Big\{\Gamma^{-}(Q^2,m_t^2)[1+a_s^2(m_t^2)A_{qq,t}^{N\!S,(2)}]\times\\
\\
& & \displaystyle \Gamma^{-}(m_t^2,m_b^2)[1+a_s^2(m_b^2)A_{qq,b}^{N\!S,(2)}]\Gamma^{v}(m_b^2,m_c^2)\\ 
\\
& & \displaystyle [1+a_s^2(m_c^2)A_{qq,c}^{N\!S,(2)}]\Gamma^{v}(m_c^2,Q_0^2)\Big\}V^{(3)}(Q^2_0)
\end{array}
\end{equation}
\item $\mathbf{V_{35}}$:
\begin{equation}
\begin{array}{rcl}
V^{(6)}_{15}(Q^2)&=&\displaystyle \Big\{\Gamma^{-}(Q^2,m_t^2)[1+a_s^2(m_t^2)A_{qq,t}^{N\!S,(2)}]\times\\
\\
& & \displaystyle \Gamma^{v}(m_t^2,m_b^2)[1+a_s^2(m_b^2)A_{qq,b}^{N\!S,(2)}]\Gamma^{v}(m_b^2,m_c^2)\\ 
\\
& & \displaystyle [1+a_s^2(m_c^2)A_{qq,c}^{N\!S,(2)}]\Gamma^{v}(m_c^2,Q_0^2)\Big\}V^{(3)}(Q^2_0)
\end{array}
\end{equation}
\item $\mathbf{T_3}$ and $\mathbf{T_8}$:
\begin{equation}
\begin{array}{rcl}
T^{(6)}_{3,8}(Q^2)&=&\displaystyle \Big\{\Gamma^{+}(Q^2,m_t^2)[1+a_s^2(m_t^2)A_{qq,t}^{N\!S,(2)}]\times\\
\\
& & \displaystyle \Gamma^{+}(m_t^2,m_b^2)[1+a_s^2(m_b^2)A_{qq,b}^{N\!S,(2)}]\Gamma^{+}(m_b^2,m_c^2)\\
\\
& & \displaystyle [1+a_s^2(m_c^2)A_{qq,c}^{N\!S,(2)}]\Gamma^{+}(m_c^2,Q_0^2)\Big\}T^{(3)}_{3,8}(Q^2_0)
\end{array}
\end{equation}

\item $\mathbf{T_{15}}$:
\begin{equation}
\begin{array}{rcl}
T^{(6)}_{15}(Q^2)&=&\displaystyle \Big\{\Gamma^{+}(Q^2,m_t^2)[1+a_s^2(m_t^2)A_{qq,t}^{N\!S,(2)}]\times\\
\\
& & \displaystyle \Gamma^{+}(m_t^2,m_b^2)[1+a_s^2(m_b^2)A_{qq,b}^{N\!S,(2)}]\Gamma^{+}(m_b^2,m_c^2)\\
\\
& & \displaystyle \begin{pmatrix} 1+a_s^2(m_c^2)[A_{qq,c}^{N\!S,(2)}-3\tilde{A}^{S,(2)}_{cq}] & -3a_s^2(m_c^2)\tilde{A}^{S,(2)}_{cg}\end{pmatrix}\underbrace{\begin{pmatrix}\Gamma_{qq}& \Gamma_{qg} \\ \Gamma_{gq} & \Gamma_{gg}\end{pmatrix}}_{(m_c^2,Q_0^2)}\Bigg\}{\Sigma^{(3)}(Q_0^2) \choose g^{(3)}(Q_0^2)}
\end{array}
\end{equation}
\item $\mathbf{T_{24}}$:
\begin{equation}
\begin{array}{rcl}
T^{(6)}_{24}(Q^2)&=&\displaystyle \Big\{\Gamma^{+}(Q^2,m_t^2)[1+a_s^2(m_t^2)A_{qq,t}^{N\!S,(2)}]\Gamma^{+}(m_t^2,m_b^2)\times\\
\\
& & \displaystyle \begin{pmatrix} 1+a_s^2(m_b^2)[A_{qq,b}^{N\!S,(2)}-4\tilde{A}^{S,(2)}_{bq}] & -4a_s^2(m_b^2)\tilde{A}^{S,(2)}_{bg}\end{pmatrix}\underbrace{\begin{pmatrix}\Gamma_{qq}& \Gamma_{qg} \\ \Gamma_{gq} & \Gamma_{gg}\end{pmatrix}}_{(m_b^2,m_c^2)}\times\\
\\
 & &\displaystyle\begin{pmatrix} M_{11}^c & M_{12}^c \\ M_{21}^c & M_{22}^c\end{pmatrix}\underbrace{\begin{pmatrix} \Gamma_{qq}& \Gamma_{qg} \\ \Gamma_{gq}& \Gamma_{gg}\end{pmatrix}}_{(m_c^2,Q_0^2)}\Bigg\}{\Sigma^{(3)}(Q_0^2) \choose g^{(3)}(Q_0^2)}
\end{array}
\end{equation}
\item $\mathbf{T_{35}}$:
\begin{equation}
\begin{array}{rcl}
T^{(6)}_{35}(Q^2)&=&\displaystyle \Big\{\Gamma^{+}(Q^2,m_t^2)\begin{pmatrix} 1+a_s^2(m_b^2)[A_{qq,t}^{N\!S,(2)}-5\tilde{A}^{S,(2)}_{tq}] & -5a_s^2(m_t^2)\tilde{A}^{S,(2)}_{tg}\end{pmatrix}\times\\
\\
& &\displaystyle\begin{pmatrix} M_{11}^b & M_{12}^b \\ M_{21}^b & M_{22}^b\end{pmatrix}\underbrace{\begin{pmatrix} \Gamma_{qq}& \Gamma_{qg} \\ \Gamma_{gq}& \Gamma_{gg}\end{pmatrix}}_{(m_b^2,m_c^2)}\\
\\
& &\displaystyle\begin{pmatrix} M_{11}^c & M_{12}^c \\ M_{21}^c & M_{22}^c\end{pmatrix}\underbrace{\begin{pmatrix} \Gamma_{qq}& \Gamma_{qg} \\ \Gamma_{gq}& \Gamma_{gg}\end{pmatrix}}_{(m_c^2,Q_0^2)}\Bigg\}{\Sigma^{(3)}(Q_0^2) \choose g^{(3)}(Q_0^2)}
\end{array}
\end{equation}
\end{itemize}

\section{Evolution and matching with intrinsic distributions}

The (quite outdated) discussion presented above concerning the PDF
matching conditions is restricted to forward evolution (\textit{i.e.}
$Q>Q_0$) and under the assumption that possible intrinsic heavy-quark
contributions are absent.\footnote{In fact, the discussion also
  assumes that the PDF matching takes place at the heavy quark-mass
  value. This is also a non-necessary assumption that we will however
  not discuss here.} This corresponds to assuming that heavy-quark
PDFs are entirely dynamically generated at the respective threshold
and are identically vanishing below threshold. While this turns out to
be a good approximation for PDFs, this is not the case for
fragmentation functions (FFs). Indeed, intrinsic heavy-quark FFs are
sizeable. In addition, assuming that there are no heavy-quark PDFs
makes backward evolution problematic. The basic problem is that the
matching-function matrix is non-squared which forbids its inversion
thus complicating backward evolution.

The purpose of this section is the relaxation of the no-intrinsic-PDF
assumption. This will imply the introduction of new matching functions
due to the presence of heavy-quark PDFs below threshold. Importantly,
the perturbative expansion of these additional functions is different
from zero at $\mathcal{O}(\alpha_s)$ even choosing $\mu_h=m_h$, with
$\mu_h$ the matching scale, producing a discontinuity already at NLO.
However, in what follows we will assume that matching conditions are
not computed beyond $\mathcal{O}(\alpha_s^2)$, \textit{i.e.}  NNLO,
and that intrinsic heavy-quark and heavy-antiquark contributions are
equal below threshold. These assumptions lead to:
\begin{equation}
h^{(n_f)}(\mu_h) = \overline{h}^{(n_f)}(\mu_h)\quad\mbox{and}\quad h^{(n_f+1)}(\mu_h) = \overline{h}^{(n_f+1)}(\mu_h)\,,
\end{equation}
where we dropped the non-scale dependence. It immediately follows that
the valence distributions $\{V,V_3,V_8,V_{15},V_{24},V_{35}\}$ match
multiplicatively as:\footnote{Note that ``1'' in
  Eq.~(\ref{eq:nsingletmc}) is to be understood as $\delta(1-x)$ in
  $x$ space.}
\begin{equation}
V_i^{(n_f+1)}(\mu_h) = [1+K_{ll}](\mu_h) V^{(n_f)}(\mu_h) =
[1+K_{ll}(\mu_h)] V_i^{(n_f)}(\mu_h)\,,\quad i = 3,8,15,24,35\,,
\label{eq:nsingletmc}
\end{equation}
with $K_{ll}$ some given function that starts at
$\mathcal{O}(\alpha_s^2)$. This equivalently amounts to saying that no
non-perturbative quark-antiquark asymmetry is originally present nor
is generated by the matching procedure.\footnote{The assumption of
  absence of non-perturbative quark-antiquark asymmetry is badly
  violated for light-valence distributions, such as up and down in the
  proton. However, we will never realistically need to match PDFs or
  FFs at the light-quark scales. On the other hand, this assumption
  helps us keep the discussion general.}

This allows us to concentrate the discussion on the singlet sector. We
start by generalising Eqs.~(\ref{eq:lightqmc}), (\ref{gluon}),
and~(\ref{eq:heavyqmc}) (changing and simplifying a bit the notation)
as:
\begin{equation}
\begin{array}{rcl}
  \displaystyle g^{(n_f+1)}&=&\displaystyle
                               [1+K_{gg}]g^{(n_f)}+K_{gl}\sum_l
                               l^{+(n_f)} + K_{gh}h^{+(n_f)}\,,\\
  \\
  \displaystyle l^{+(n_f+1)}&=&\displaystyle [1+K_{ll}]l^{+(n_f)}\,,\\
  \\
  \displaystyle h^{+(n_f+1)}&=&\displaystyle K_{hg}g^{(n_f)} + K_{hl}\sum_l
                                l^{+(n_f)}+[1+K_{hh}]h^{+(n_f)}\,,\\
  \\
  H^{+(n_f+1)} &=& H^{+(n_f)}\,,
\end{array}
\label{eq:IntMatchingConds}
\end{equation}
where $q^+=q+\overline{q}$ and $l$ runs between $1$ and $n_f$. We
remind that $h$ is the $(n_f+1)$-th heavier flavour, \textit{i.e.}
the one that becomes active at the threshold under consideration. We
have also introduced the ``super'' heavy quark flavour $H$, heavier
than $h$, that is not affected by the matching at the $n_f$-th
threshold. Note that there are two new functions coming into play,
namely $K_{gh}$ and $K_{hh}$, both multiplying the intrinsic
heavy-quark contribution $h^{+(n_f)}$. Moreover, we have written
Eq.~(\ref{eq:IntMatchingConds}) in such a way that the coefficients
$K$ start at least at $\mathcal{O}(\alpha_s)$. In fact, we have
already seen that $K_{ll}=\mathcal{O}(\alpha_s^2)$; the same applies
to $K_{hl}$ and $K_{gl}$. Defining the column vector $P$ with
components $(g,d^+,u^+,s^+,c^+,b^+,t^+)$, we can represent the
matching in a matrix form as:
\begin{equation}
P^{(n_f+1)} = \left[\mathbb{I}+\mathbb{K}^{(n_f)}\right]P^{(n_f)}\,,
\end{equation}
where $\mathbb{I}$ is the $7\times7$ unity matrix and
$\mathbb{K}^{(n_f)}$ are given, depending on the value of $n_f$, by:
\begin{equation}
\begin{array}{ll}
\mathbb{K}^{(0)} = 
\begin{pmatrix}
K_{gg} & K_{gh} & 0  & 0  & 0  & 0  & 0 \\
K_{hg} & K_{hh} & 0  & 0  & 0  & 0  & 0 \\
0 & 0 & 0  & 0  & 0  & 0  & 0 \\
0 & 0 & 0  & 0  & 0  & 0  & 0 \\
0 & 0 & 0  & 0  & 0  & 0  & 0 \\
0 & 0 & 0  & 0  & 0  & 0  & 0 \\
0 & 0 & 0  & 0  & 0  & 0  & 0
\end{pmatrix}
&
\mathbb{K}^{(1)} = 
\begin{pmatrix}
K_{gg} & K_{gl} & K_{gh} & 0  & 0  & 0  & 0 \\
0 & K_{ll} & 0 & 0  & 0  & 0  & 0  \\
K_{hg} & K_{hl} & K_{hh}  & 0  & 0  & 0  & 0 \\
0 & 0 & 0  & 0  & 0  & 0  & 0 \\
0 & 0 & 0  & 0  & 0  & 0  & 0 \\
0 & 0 & 0  & 0  & 0  & 0  & 0 \\
0 & 0 & 0  & 0  & 0  & 0  & 0
\end{pmatrix}
\\ \\
\mathbb{K}^{(2)} = 
\begin{pmatrix}
K_{gg} & K_{gl} & K_{gl} & K_{gh} & 0  & 0  & 0 \\
0 & K_{ll} & 0 & 0  & 0  & 0  & 0  \\
0 & 0 & K_{ll} & 0  & 0  & 0  & 0  \\
K_{hg} & K_{hl} & K_{hl} & K_{hh}  & 0  & 0  & 0 \\
0 & 0 & 0  & 0  & 0  & 0  & 0 \\
0 & 0 & 0  & 0  & 0  & 0  & 0 \\
0 & 0 & 0  & 0  & 0  & 0  & 0
\end{pmatrix}
&
\mathbb{K}^{(3)} = 
\begin{pmatrix}
K_{gg} & K_{gl} & K_{gl} & K_{gl} & K_{gh}  & 0  & 0 \\
0 & K_{ll} & 0 & 0  & 0  & 0  & 0  \\
0 & 0 & K_{ll} & 0  & 0  & 0  & 0  \\
0 & 0 & 0 & K_{ll} & 0  & 0  & 0  \\
K_{hg} & K_{hl} & K_{hl} & K_{hl} & K_{hh}  & 0  & 0 \\
0 & 0 & 0  & 0  & 0  & 0  & 0 \\
0 & 0 & 0  & 0  & 0  & 0  & 0
\end{pmatrix}
\\ \\
\mathbb{K}^{(4)} = 
\begin{pmatrix}
K_{gg} & K_{gl} & K_{gl} & K_{gl} & K_{gl} & K_{gh} & 0  \\
0 & K_{ll} & 0 & 0  & 0  & 0  & 0  \\
0 & 0 & K_{ll} & 0  & 0  & 0  & 0  \\
0 & 0 & 0 & K_{ll} & 0  & 0  & 0  \\
0 & 0 & 0 & 0 & K_{ll}  & 0  & 0  \\
K_{hg} & K_{hl} & K_{hl} & K_{hl} & K_{hl} & K_{hh}  & 0 \\
0 & 0 & 0  & 0  & 0  & 0  & 0
\end{pmatrix}
&
\mathbb{K}^{(5)} = 
\begin{pmatrix}
K_{gg} & K_{gl} & K_{gl} & K_{gl} & K_{gl} & K_{gl} & K_{gh} \\
0 & K_{ll} & 0 & 0 & 0  & 0  & 0  \\
0 & 0 & K_{ll} & 0 & 0  & 0  & 0  \\
0 & 0 & 0 & K_{ll} & 0  & 0  & 0  \\
0 & 0 & 0 & 0 & K_{ll} & 0  & 0  \\
0 & 0 & 0 & 0 & 0 & K_{ll} & 0  \\
K_{hg} & K_{hl} & K_{hl} & K_{hl} & K_{hl} & K_{hl} & K_{hh} \\
\end{pmatrix}
\end{array}
\label{eq:MathcingMatrices}
\end{equation}
The algorithm to determine the matrix $\mathbb{K}^{(n_f)}$, indexing
the gluon distribution with 0 and the quark distributions from 1 to 6,
reads:
\begin{equation}
\mathbb{K}_{ij}^{(n_f)} =
\left\{
\begin{array}{ll}
K_{gg} & \quad i = j = 0\\
K_{gl} & \quad i = 0\,,\;1\leq j \leq n_f\\
K_{gh} & \quad i = 0\,,\;j = n_f+1\\
K_{ll} & \quad i = j\,,\;1\leq i,j \leq n_f\\
K_{hg} & \quad i = n_f+1\,,\;j = 0\\
K_{hl} & \quad i = n_f+1\,,\;1\leq j \leq n_f\\
K_{hh} & \quad i = j = n_f+1\\
0 & \quad\mbox{elsewhere}
\end{array}
\right.\,.
\label{eq:algorithm}
\end{equation}

Now, in order to apply the matching condition to the distributions in
the evolution basis $E=(g,\Sigma,-T_3,T_8,T_{15},T_{24},T_{35})$
(notice the minus sign in front of $T_3$), it is necessary to rotate
the matching matrices from the physical basis to the evolution basis
through $T\mathbb{K}^{(n_f)}T^{-1}$, where $T$ transforms the vector
$P$ into the evolution-basis vector $E$:
\begin{equation}
T = 
\begin{pmatrix}
1 & 0 & 0  & 0  & 0  & 0  & 0 \\
0 & 1 & 1  & 1  & 1  & 1  & 1 \\
0 & 1 & -1  & 0  & 0  & 0  & 0 \\
0 & 1 & 1  & -2  & 0  & 0  & 0 \\
0 & 1 & 1  & 1  & -3  & 0  & 0 \\
0 & 1 & 1  & 1  & 1  & -4  & 0 \\
0 & 1 & 1  & 1  & 1  & 1  & -5
\end{pmatrix}\,.
\label{eq:RotationMatrix}
\end{equation}
The resulting matching matrices are much more complicated and
generally less sparse than those in
Eq.~(\ref{eq:MathcingMatrices}). However, we can still compute them
algorithmically depending on the number of light flavours $n_f$. To do
so, we define the following linearly independent combinations:
\begin{equation}
\begin{array}{lcl}
M_1 &=& K_{gg}\,,\\
M_2 &=& K_{gh} + n_f K_{gl}\,,\\
M_3 &=& K_{gh}-K_{gl}\,,\\
M_4 &=& K_{hg}\,,\\
M_5 &=& K_{hh}+n_f(K_{hl}+K_{ll}) \,,\\
M_6 &=& K_{hh}-(K_{hl}+K_{ll})\,,\\
M_7 &=& K_{ll}\,,
\end{array}
\label{eq:Mcombs}
\end{equation}
such that:
\begin{equation}
\footnotesize
\left[T\mathbb{K}^{(n_f)}T^{-1}\right]_{ij} =
\left\{
\begin{array}{lll}
M_1 & \quad i = 0 &\quad j = 0\\
\frac{1}{6}M_2 & \quad i = 0 &\quad  j = 1\\
0 & \quad i = 0 &\quad  2\leq j \leq n_f\\
-\frac{1}{n_f+1}M_3 & \quad i = 0 &\quad  j = n_f+1\\
\frac{1}{j(j-1)}M_2 & \quad i = 0 &\quad  n_f+ 2\leq j \leq 6\\
\\
M_4 & \quad i = 1&\quad  j=0\\
\frac16 M_5 & \quad i = 1 &\quad  j=1\\
0 & \quad i = 1 &\quad  2\leq j \leq n_f\\
-\frac{1}{n_f+1}M_6 & \quad i = 1 &\quad  j = n_f+1\\
\frac{1}{j(j-1)}M_5 & \quad i = 1 &\quad  n_f+
                                                            2\leq j
                                                            \leq 6\\
\\
\delta_{ij}M_7&\quad 2\leq i \leq n_f &\quad 0\leq j
                                                         \leq 6\\
\\
-n_fM_4 &\quad i = n_f+1 &\quad j =0\\
\frac{n_f}6\left(-M_5+(n_f+1)M_7\right) &\quad i = n_f+1 &\quad j =1\\
0 & \quad i = n_f+1 &\quad  2\leq j \leq n_f\\
\frac1{n_f+1}\left(n_f M_6+(n_f+1)M_7\right) &\quad i = n_f+1 &\quad j =n_f+1\\
\frac{n_f}{j(j-1)}\left( -M_5+(n_f+1)M_7\right) &\quad i = n_f+1
                  &\quad n_f+2 \leq j \leq 6\\
\\
M_4 & \quad n_f+2\leq i \leq 6 &\quad  j=0\\
\frac16 M_5 & \quad n_f+2\leq i \leq 6 &\quad  j=1\\
0 & \quad n_f+2\leq i \leq 6 &\quad  2\leq j \leq n_f\\
-\frac{1}{n_f+1}M_6 & \quad n_f+2\leq i \leq 6 &\quad  j = n_f+1\\
\frac{1}{j(j-1)}M_5 & \quad n_f+2\leq i \leq 6 &\quad  n_f+
                                                            2\leq j
                                                            \leq 6\\
\end{array}
\right.
\label{eq:algorithmEv}
\end{equation}

In order to be able to perform the backward evolution, it is necessary
to invert the matching matrix $\mathbb{I}+\mathbb{K}^{(n_f)}$. This is
conveniently done perturbatively by expanding the matrix
$\left[\mathbb{I}+\mathbb{K}^{(n_f)}\right]^{-1}$ in powers of
$\alpha_s$ and truncating the expansion at the appropriate
order. Given the assumption discussed above, we are not interested in
going beyond $\mathcal{O}(\alpha_s^2)$. In addition, we know that:
\begin{equation}
\mathbb{K}^{(n_f)} = a_s \mathbb{K}^{(1)(n_f)}+a_s^2 \mathbb{K}^{(2)(n_f)}+\mathcal{O}(\alpha_s^3)\,,
\end{equation}
so that:
\begin{equation}
\left[\mathbb{I}+\mathbb{K}^{(n_f)}\right]^{-1} = 1 - a_s
\mathbb{K}^{(1)(n_f)}-a_s^2\left[\mathbb{K}^{(2)(n_f)}-\left(\mathbb{K}^{(1)(n_f)}\right)^2\right]+\mathcal{O}(\alpha_s^3)\,.
\label{eq:inversion}
\end{equation}
But the matrix $\mathbb{K}^{(1)(n_f)}$ is particularly simple because
at $\mathcal{O}(\alpha_s)$ the matching functions $K_{ll}$, $K_{hl}$,
and $K_{gl}$ vanish. Therefore, using Eq.~(\ref{eq:algorithm}) we
find:
\begin{equation}
\mathbb{K}_{ij}^{(1)(n_f)} =\delta_{i,0}\left(\delta_{j,0}K_{gg}^{(1)} +
\delta_{j,n_f+1}K_{gh}^{(1)}\right)+ \delta_{i,n_f+1}\left(\delta_{j,0}K_{hg}^{(1)}
+ \delta_{j,n_f+1}K_{hh}^{(1)}\right)\,,
\end{equation}
such that its square has the same structure and can be written as:
\begin{equation}
\begin{array}{rcl}
\left(\mathbb{K}^{(1)(n_f)}\right) _{ij}^2 &=&\displaystyle \delta_{i,0}\left(\delta_{j,0}\sum_{k=g,h}K_{gk}^{(1)} K_{kg}^{(1)} +
\delta_{j,n_f+1}\sum_{k=g,h}K_{gk}^{(1)} K_{kh}^{(1)}\right)\\
\\
  &+&\displaystyle  \delta_{i,n_f+1}\left(\delta_{j,0}\sum_{k=g,h}K_{hk}^{(1)} K_{kg}^{(1)}
+ \delta_{j,n_f+1}\sum_{k=g,h}K_{hk}^{(1)} K_{kh}^{(1)}\right)\,,
\end{array}
\label{eq:squareK1}
\end{equation}
where, in $x$ space, the product of two $K$ functions is to be
understood as a convolution. We report below the expression in $x$
space of the functions $K_{gh}^{(1)}$ and $K_{gh}^{(1)}$ for the
matching of PDFs:
\begin{equation}
\begin{array}{rcl}
\displaystyle K_{gh}^{(1)}(x) &=&\displaystyle 
2C_F\frac{1+(1-x)^2}{x}\left(\ln\frac{\mu_h^2}{m_h^2}-1-2\ln x\right)\,,\\
\\
\displaystyle K_{hh}^{(1)}(x) &=&\displaystyle 
2C_F\left[\frac{1+x^2}{1-x}\left(\ln\frac{\mu_h^2}{m_h^2}-1-2\ln(1-x)\right)\right]_+\,.
\end{array}
\end{equation}
We notice that the $\mathcal{O}(\alpha_s^2)$ correction to these
functions is currently unknown, therefore we only use the
$\mathcal{O}(\alpha_s)$ expressions. For the sake of completeness, we
also report here the remaining $\mathcal{O}(\alpha_s)$ matching
functions:
\begin{equation}
\begin{array}{rcl}
\displaystyle K_{gg}^{(1)}(x) &=&\displaystyle -\frac{4}{3}T_R\delta(1-x)\ln\frac{\mu_h^2}{m_h^2}\,,\\
\\
\displaystyle K_{hg}^{(1)}(x) &=&\displaystyle 4T_R\left[x^2+(1-x)^2\right]\ln\frac{\mu_h^2}{m_h^2}\,,
\end{array}
\end{equation}
Of course, the matrix
$\left[\mathbb{I}+\mathbb{K}^{(n_f)}\right]^{-1}$ in
Eq.~(\ref{eq:inversion}) has to be rotated into the evolution basis by
means of the transformation matrix $T$ in
Eq.~(\ref{eq:RotationMatrix})

In order to implement the possibility to compute the evolution
operator in the VFNS, we need to know how the evolution operator for a
given number of active flavours $n_f$ matches at the following
threshold. The first step is to generalise the structure of the DGLAP
evolution equations in the QCD evolution basis allowing for the
presence of (scale-independent) intrinsic contributions. It is
important to realise that the DGLAP equations govern the evolution in
$\mu$ of the collinear distributions but do not give us any
information on the non-dynamic part of these distributions,
\textit{i.e.} its intrinsic contribution. The starting point is the
DGLAP equation for the quark distribution $q_k^+$\footnote{We stick to
  the assumption that intrinsic contributions are symmetric upon
  charge conjugation. This allow us to forget about the valence sector
  that will always evolve multiplicatively.} in its basic form:
\begin{equation}
\frac{dq_k^+}{d\ln\mu^2} =\frac{dq_k^+}{dt}= 2P_{qg}\otimes g+
(P_{qq}^V+P_{q\overline{q}}^V)\otimes
q_k^++(P_{qq}^S+P_{q\overline{q}}^S)\otimes \sum_{i=1}^{n_f} q_i^+\,.
\end{equation}
It is crucial to notice that the sum on the r.h.s. of the equation
above runs over the $n_f$ active flavours only. Now, we sum the index
$k$ up to $n_f$ and we get:
\begin{equation}
\frac{d}{dt}\sum_{k=1}^{n_f} q_k^+= 2n_f P_{qg}\otimes g+
P_{qq}\otimes \sum_{k=1}^{n_f} q_k^+\,,
\end{equation}
where we have defined:
\begin{equation}
P_{qq} = (P_{qq}^V+P_{q\overline{q}}^V)+n_f(P_{qq}^S+P_{q\overline{q}}^S).
\end{equation}
In the presence of intrinsic heavy-quark distributions and after some
algebra, one finds that:
\begin{equation}
\frac{d\Sigma}{dt}= 2n_f P_{qg}\otimes g+n_fP_{qq}\otimes
\left(\frac{\Sigma}{6}+\sum_{j=n_f+1}^6\frac{E_j}{j(j-1)}\right)\,,
\label{eq:singletDGLAP}
\end{equation}
with $E_j\in \{\Sigma,-T_3,T_8,T_{15},T_{24},T_{35}\}$. Therefore, the
singlet distribution, on top of the gluon, couples also to the other
quark distributions. Notice that:
\begin{equation}
\sum_{j=n_f+1}^6\frac{1}{j(j-1)} = \frac{1}{n_f}-\frac16\,,
\end{equation}
such that, for $E_j = \Sigma$, Eq.~(\ref{eq:singletDGLAP}) reproduces
the usual form of the DGLAP equation for the singlet. We now need to
determine how the distributions $E_j$, for $j>1$,
evolve. We find:
\begin{equation}
  \frac{dE_j}{dt}= (1-\theta_{n_fj})\left[2n_f P_{qg}\otimes
    g+n_fP_{qq}\otimes
    \left(\frac{\Sigma}{6}+\sum_{i=n_f+1}^6\frac{E_i}{i(i-1)}\right)\right]+\theta_{n_fj}P^+\otimes
  E_j\,,
\label{eq:nonsingletDGLAP}
\end{equation}
with:
\begin{equation}
\theta_{n_fj}=\left\{
\begin{array}{ll}
1   & \quad n_f\geq j\\
0 & \quad n_f < j
\end{array}
\right.\,. 
\end{equation}
Therefore, as expected, for a given $n_f$ (and thus a given energy)
the distribution $E_j$ evolves multiplicatively through $P^+$ if
$j\leq n_f$ and evolves exactly like the singlet if $j>n_f$. As an
example, for $n_f=4$, that is to say for energies between the charm
and the bottom thresholds, $E_2=T_3$, $E_3=T_8$, and $E_4=T_{15}$
evolve multiplicatively, while $E_5=T_{24}$ and $E_6=T_{35}$ evolve
like the singlet.

Finally, the gluon distribution evolves as:
\begin{equation}
\frac{dg}{dt}= P_{gg}\otimes g+
P_{gq}\otimes \sum_{i=1}^{n_f} q_i^+\,,
\end{equation}
that translates into:
\begin{equation}
\frac{dg}{dt}= P_{gg}\otimes g+n_fP_{gq}\otimes
\left(\frac{\Sigma}{6}+\sum_{j=n_f+1}^6\frac{E_j}{j(j-1)}\right)\,.
\label{eq:gluonDGLAP}
\end{equation}
Eqs.~(\ref{eq:singletDGLAP}), (\ref{eq:nonsingletDGLAP}), and
(\ref{eq:gluonDGLAP}) fully define the splitting kernel matrix in the
general case of possible presence of intrinsic heavy-quark
contributions. Noticeably, these equations correctly reduce to the
more familiar ones if these contributions are set to zero.

Therefore, the evolution of the vector of distributions $E$ with $n_f$
active flavours and possible intrinsic heavy-quark contributions takes
the matricial form:
\begin{equation}
\frac{dE}{dt} = \mathbb{P}^{(n_f)}\otimes E\,,
\label{eq:DGLAPE}
\end{equation}
where the entries of the splitting-function matrix
$\mathbb{P}^{(n_f)}$ can be defined algorithmically as:
\begin{equation}
\mathbb{P}_{ij}^{(n_f)}=
\left\{
\begin{array}{lll}
P_{gg} & \quad i = 0 &\quad j = 0\\
\frac{n_f}{6}P_{gq} & \quad i = 0 & \quad j=1\\
0 &\quad i = 0 & \quad 2 \leq j \leq n_f\\
\frac{6}{j(j-1)}\frac{n_f}{6}P_{gq} &\quad i = 0 & \quad n_f+1 \leq j \leq 6\\
\\
2n_fP_{qg} & \quad i = 1 &\quad j = 0\\
\frac{n_f}{6}P_{qq} & \quad i = 1 & \quad j=1\\
0 &\quad i = 1 & \quad 2 \leq j \leq n_f\\
\frac{6}{j(j-1)}\frac{n_f}{6}P_{qq} &\quad i = 1 & \quad n_f+1 \leq j \leq 6\\
\\
\delta_{ij}P^+ & \quad 2 \leq i \leq n_f & \quad 0 \leq j \leq 6\\
\\
2n_fP_{qg} & \quad n_f+1 \leq i \leq 6  &\quad j = 0\\
\frac{n_f}{6}P_{qq} & \quad n_f+1 \leq i \leq 6  & \quad j=1\\
0 &\quad n_f+1 \leq i \leq 6 & \quad 2 \leq j \leq n_f\\
\frac{6}{j(j-1)}\frac{n_f}{6}P_{qq} &\quad n_f+1 \leq i \leq 6 & \quad n_f+1 \leq j \leq 6\\
\end{array}
\right.\,.
\label{eq:splittingalg}
\end{equation}

\subsection{Evolution operator}

We assume that the set of distributions $E$ evolves through the
operator ${\bm\Gamma}$ as:
\begin{equation}
  E(t) = {\bm\Gamma}(t,t_0)\otimes E(t_0)\,.
\end{equation}
In order to achieve an efficient computation of the evolution
operator, we need to identify the general structure of ${\bm \Gamma}$.
Given the possible presence of heavy-quark thresholds in the evolution
interval $[t_0,t]$, the operator ${\bm\Gamma}$ is in fact a product of
operators. More specifically, it can be written as:
\begin{equation}
{\bm\Gamma}(t,t_0)= {\bm\Gamma}^{(N)}(t,t_N)
\otimes\prod_{n_f=N-1}^{0}
\mathcal{M}^{(n_f)}\otimes{\bm\Gamma}^{(n_f)}(t_{n_f+1}-\epsilon,t_{n_f})
\label{eq:conbevop}
\end{equation}
with
$\mathcal{M}^{(n_f)} =
T\left[\mathbb{I}+\mathbb{K}^{(n_f)}\right]T^{-1}$
given in Eq.~(\ref{eq:algorithmEv}) and $t>t_N>\dots>t_1>t_0$, being
$t_1,\dots,t_N$ the heavy-quark thresholds enclosed in the evolution
interval. Using Eq.~(\ref{eq:DGLAPE}), each evolution operator
${\bm\Gamma}^{(n_f)}$ obeys the evolution equation:
\begin{equation}
  \frac{d}{dt}{\bm\Gamma}^{(n_f)}(t,t_{n_f}) =
  \mathbb{P}^{(n_f)}(t)\otimes {\bm\Gamma}^{(n_f)}(t,t_{n_f})\,,
\label{eq:DGLAPG}
\end{equation}
with boundary condition
${\bm\Gamma}^{(n_f)}(t_{n_f},t_{n_f})=\mathbb{I}$. The solution of the
equation above can then be written as:
\begin{equation}
  {\bm\Gamma}^{(n_f)}(t,t_{n_f}) = \mathcal{P}\exp\left[\int_{t_{n_f}}^tdt'\,\mathbb{P}^{(n_f)}(t')\right]\,,
\end{equation}
where $\mathcal{P}$ symbolises the path ordering. Given the structure
of the splitting-function matrix in Eq.~(\ref{eq:splittingalg}), one
can infer the structure of ${\bm \Gamma}^{(n_f)}$. It turns out that
the structure of ${\bm \Gamma^{(n_f)}}$ is almost the same as that of
$\mathbb{P}^{(n_f)}$. Specifically:
\begin{equation}
  {\bm\Gamma}^{(n_f)}_{ij}=\mathbb{I}_{ij}+
  \left\{
\begin{array}{lll}
\Gamma_{gg} & \quad i = 0 &\quad j = 0\\
\Gamma_{gq} & \quad i = 0 & \quad j=1\\
0 &\quad i = 0 & \quad 2 \leq j \leq n_f\\
\frac{6}{j(j-1)}\Gamma_{gq} &\quad i = 0 & \quad n_f+1 \leq j \leq 6\\
\\
\Gamma_{qg} & \quad i = 1 & \quad j = 0\\
\Gamma_{qq} & \quad i = 1 & \quad j = 1\\
0 &\quad i = 1 & \quad 2 \leq j \leq n_f\\
\frac{6}{j(j-1)}\Gamma_{qq} &\quad i = 1 & \quad n_f+1 \leq j \leq 6\\
\\
\delta_{ij}\Gamma^+ & \quad 2 \leq i \leq n_f & \quad 0 \leq j \leq 6\\
\\
\Gamma_{qg} & \quad n_f+1 \leq i \leq 6  &\quad j = 0\\
\Gamma_{qq} & \quad n_f+1 \leq i \leq 6  & \quad j=1\\
0 &\quad n_f+1 \leq i \leq 6 & \quad 2 \leq j \leq n_f\\
\frac{6}{j(j-1)}\Gamma_{qq} &\quad n_f+1 \leq i \leq 6 & \quad n_f+1 \leq j \leq 6\\
\end{array}
\right.\,.
\label{eq:evolopalg}
\end{equation}
% Using Eqs.~(\ref{eq:splittingalg}) and~(\ref{eq:evolopalg}), we can
% work out the product $\mathbb{P}^{(n_f)}{\bm\Gamma}^{(n_f)}$ required
% for the computation of the evolution operator. This reads:
% \begin{equation}\label{eq:PGamma}
% \left[\mathbb{P}^{(n_f)}{\bm\Gamma}^{(n_f)}\right]_{ij}=\mathbb{P}^{(n_f)}_{ij}+
% \left\{
% \begin{array}{lll}
% P_{gg}\Gamma_{gg}+P_{gq}\Gamma_{qg} & \quad i = 0 &\quad j = 0\\
% P_{gg}\Gamma_{gq}+P_{gq}\Gamma_{qq} & \quad i = 0 & \quad j=1\\
% 0 &\quad i = 0 & \quad 2 \leq j \leq n_f\\
% \frac{6}{j(j-1)}\left[ P_{gg}\Gamma_{gq}+P_{gq}\Gamma_{qq}\right]&\quad i = 0 & \quad n_f+1 \leq j \leq 6\\
% \\
% 2n_fP_{qg}\Gamma_{gg}+P_{qq}\Gamma_{qg} & \quad i = 1 &\quad j = 0\\
% 2n_fP_{qg}\Gamma_{gq}+P_{qq}\Gamma_{qq} & \quad i = 1 & \quad j=1\\
% 0 &\quad i = 1 & \quad 2 \leq j \leq n_f\\
% \frac{6}{j(j-1)} \left[2n_fP_{qg}\Gamma_{gq}+P_{qq}\Gamma_{qq}\right]&\quad i = 1 & \quad n_f+1 \leq j \leq 6\\
% \\
% \delta_{ij}P^+\Gamma^+ & \quad 2 \leq i \leq n_f & \quad 0 \leq j \leq 6\\
% \\
% \frac{6}{j(j-1)}\left[2n_fP_{qg}\Gamma_{gq}+P_{qq}\Gamma_{qq} \right]&\quad n_f+1 \leq i \leq 6 & \quad n_f+1 \leq j \leq 6\\
% \end{array}
% \right.\,.
% \end{equation}

The next step is the computation of the product
$\mathcal{M}^{(n_f)}{\bm\Gamma}^{(n_f)}=T\left[\mathbb{I}+\mathbb{K}^{(n_f)}\right]T^{-1}{\bm\Gamma}^{(n_f)}$
appearing in Eq.~(\ref{eq:conbevop}) that we conveniently split as:
\begin{equation}
\mathcal{M}^{(n_f)}{\bm\Gamma}^{(n_f)}={\bm \Gamma}^{(n_f)}+
T\mathbb{K}^{(n_f)}T^{-1}+\mathbb{S}^{(n_f)}\,,
\label{eq:nthfactor}
\end{equation}
where we have defined
$\mathbb{S}^{(n_f)}=T\mathbb{K}^{(n_f)}
T^{-1}\left[{\bm\Gamma}^{(n_f)}-\mathbb{I}\right]$ that is given by:
% \begin{equation}
% \footnotesize
% \mathbb{S}_{ik}^{(n_f)} =
% \left\{
% \begin{array}{lll}
% K_{gg}\Gamma_{gg}+K_{gl}\Gamma_{qg} &
%                                                                       \quad i = 0 &\quad k = 0\\
% K_{gg}\Gamma_{gq}+K_{gl}\Gamma_{qq} & \quad i = 0 &\quad k =1\\
% 0 & \quad i = 0 &\quad  2\leq k \leq n_f\\
% \frac1{k(k-1)}\left[K_{gg}\Gamma_{gq}+K_{gl}\Gamma_{qq}\right] & \quad i = 0 &\quad  n_f+ 1\leq k \leq 6\\
% \\
% K_{hg}\Gamma_{gg}+K_{hl}\Gamma_{qg}+K_{ll}\Gamma_{qg} & \quad i = 1&\quad  k=0\\
% K_{hg}\Gamma_{gq}+K_{hl}\Gamma_{qq}+K_{ll}\Gamma_{qq}& \quad i = 1 &\quad  k=1\\
% 0 & \quad i = 1 &\quad  2\leq k \leq n_f\\
% \frac{1}{k(k-1)}\left[K_{hg}\Gamma_{gq}+K_{hl}\Gamma_{qq}+K_{ll}\Gamma_{qq}\right] & \quad i = 1 &\quad  n_f+1\leq k\leq 6\\
% \\
% \delta_{ik}K_{ll}\Gamma^+&\quad 2\leq i \leq n_f &\quad 0\leq k
%                                                          \leq 6\\
% \\
% -n_f \left(K_{hg}\Gamma_{gg}+K_{hl}\Gamma_{qg}\right)+K_{ll}\Gamma_{qg}&\quad i = n_f+1 &\quad k =0\\
% -n_f \left(K_{hg}\Gamma_{gq}+K_{hl}\Gamma_{qq}\right)+K_{ll}\Gamma_{qq}&\quad i = n_f+1 &\quad k =1\\
% 0 & \quad i = n_f+1 &\quad  2\leq k \leq n_f\\
%   \frac{6}{k(k-1)}\left[-n_f \left(K_{hg}\Gamma_{gq}+K_{hl}\Gamma_{qq}\right)+K_{ll}\Gamma_{qq}\right]
% &\quad i = n_f+1 &\quad n_f+1 \leq k \leq 6\\
% \\
%   \frac{6}{k(k-1)}\left[K_{hg}\Gamma_{gq}+K_{hl}\Gamma_{qq}+K_{ll}\Gamma_{qq} \right]
% & \quad n_f+2\leq i \leq 6 &\quad  n_f+
%                                                             2\leq k
%                                                             \leq 6\\
% \end{array}
% \right.
% \end{equation}
\begin{equation}
\footnotesize
\mathbb{S}_{ij}^{(n_f)} =
\left\{
\begin{array}{lll}
M_1\Gamma_{gg}+\frac{1}{1+n_f}\left(M_2-M_3\right)\Gamma_{qg} &
                                                                      \quad i = 0 &\quad j = 0\\
M_1\Gamma_{gq}+\frac{1}{1+n_f}\left(M_2-M_3\right)\Gamma_{qq} & \quad i = 0 &\quad j =1\\
0 & \quad i = 0 &\quad  2\leq j \leq n_f\\
\frac1{j(j-1)}\left[M_1\Gamma_{gq}+\frac{1}{1+n_f}\left(M_2-M_3\right)\Gamma_{qq}\right] & \quad i = 0 &\quad  n_f+ 1\leq j \leq 6\\
\\
M_4\Gamma_{gg}+\frac{1}{n_f+1}\left(M_5-M_6\right)\Gamma_{qg} & \quad i = 1&\quad  j=0\\
M_4\Gamma_{gq}+\frac{1}{n_f+1}\left(M_5-M_6\right)\Gamma_{qq}& \quad i = 1 &\quad  j=1\\
0 & \quad i = 1 &\quad  2\leq j \leq n_f\\
\frac{1}{j(j-1)}\left[M_4\Gamma_{gq}+\frac{1}{n_f+1}(M_5-M_6)\Gamma_{qq}\right] & \quad i = 1 &\quad  n_f+1\leq j\leq 6\\
\\
\delta_{ij}M_7\Gamma^+&\quad 2\leq i \leq n_f &\quad 0\leq j
                                                         \leq 6\\
\\
-n_f\left(M_4\Gamma_{gg}+\frac{1}{n_f+1}(M_5-M_6)\Gamma_{qg}\right)+(n_f+1)M_7\Gamma_{qg}&\quad i = n_f+1 &\quad j =0\\
-n_f\left(M_4\Gamma_{gq}+\frac{1}{n_f+1}(M_5-M_6)\Gamma_{qq}\right)+(n_f+1)M_7\Gamma_{qq}&\quad i = n_f+1 &\quad j =1\\
0 & \quad i = n_f+1 &\quad  2\leq j \leq n_f\\
  \frac{6}{j(j-1)}\left[-n_f\left(M_4\Gamma_{gq}+\frac{1}{n_f+1}(M_5-M_6)\Gamma_{qq}\right)+(n_f+1)M_7\Gamma_{qq}\right]
&\quad i = n_f+1 &\quad n_f+1 \leq j \leq 6\\
\\
  \frac{6}{j(j-1)}\left[M_4\Gamma_{gq}+\frac{1}{n_f+1}(M_5-M_6)\Gamma_{qq} \right]
& \quad n_f+2\leq i \leq 6 &\quad  n_f+
                                                            2\leq j
                                                            \leq 6\\
\end{array}
\right.
\label{eq:Smatrix}
\end{equation}
Eq.~(\ref{eq:nthfactor}) along with Eq.~(\ref{eq:Smatrix}) allow us to
infer the general structure on the $n_f$-th factor in
Eq.~(\ref{eq:conbevop}). It should be clear that the structure is
driven by the term $T\mathbb{K}^{(n_f)}T^{-1}$ in
Eq.~(\ref{eq:algorithmEv}).

We finally comment on FFs. So far, to the best of my knowledge, none
of the $\mathcal{O}(\alpha_s^2)$ corrections to the FF matching
conditions has been computed. Therefore, we must limit to
$\mathbb{K}^{(1)(n_f)}$ that is simply given by the transpose of that
for PDFs.

\subsection{In the presence of N$^3$LO corrections}

If $\mathcal{O}(\alpha_s^3)$ corrections are included in the matching, matching
conditions in Eq.~(\ref{eq:IntMatchingConds}) need to be generalised as follows:
\begin{equation}
\begin{array}{rcl}
  \displaystyle g^{(n_f+1)}&=&\displaystyle
                               [1+K_{gg}]g^{(n_f)}+K_{gl}\sum_l
                               l^{+(n_f)} + K_{gh}h^{+(n_f)}\,,\\
  \\
  \displaystyle l^{+(n_f+1)}&=&\displaystyle [1+K_{ll}]l^{+(n_f)}+K_{ll'}\sum_l l^{+(n_f)}+K_{lg}g^{(n_f)}\,,\\
  \\
  \displaystyle h^{+(n_f)}&=&\displaystyle K_{hg}g^{(n_f)} + K_{hl}\sum_l
                                l^{+(n_f)}+[1+K_{hh}]h^{+(n_f)}\,,\\
  \\
  H^{+(n_f+1)} &=& H^{+(n_f)}\,,\\
\\
  \displaystyle l^{-(n_f+1)}&=&\displaystyle [1+K_{ll}^-]l^{-(n_f)}\,,\\
\end{array}
\label{eq:IntMatchingCondsN3LO}
\end{equation}
where the new matching functions $K_{ll'}$ and $K_{lg}$, which start
at $\mathcal{O}(\alpha_s^3)$, are introduced and where we neglected a
possible $K_{lh}$ that in principle exists but will not probably be
computed in the next future. Also, we introduced the matching function
$K_{ll}^-$ which differs from $K_{ll}$ starting from
$\mathcal{O}(\alpha_s^3)$. With this modification, the matching
algorithm in the evolution basis in Eq.~(\ref{eq:algorithmEv}) is more
complicated. First we need to define two more combinations
w.r.t. Eq.~(\ref{eq:Mcombs}), \textit{i.e.} $M_8$ and $M_9$:
\begin{equation}
  \begin{array}{lcl}
  M_1 &=& K_{gg}\,,\\
  M_2 &=& K_{gh} + n_f K_{gl}\,,\\
  M_3 &=& K_{gh}-K_{gl}\,,\\
  M_4 &=& K_{hg}+n_f K_{lg}\,,\\
  M_5 &=& K_{hh}+n_f(K_{hl}+K_{ll}+n_f K_{ll'}) \,,\\
  M_6 &=& K_{hh}-(K_{hl}+K_{ll}+n_f K_{ll'})\,,\\
  M_7 &=& K_{ll}\,,\\
  M_8 &=& K_{hg} -  K_{lg}\,,\\
  M_9 &=& K_{hh}+n_f(K_{hl}+K_{ll}-K_{ll'})\,,\\
  M_{10} &=& K_{ll}^-\,,
  \end{array}
  \label{eq:McombsN3LO}
\end{equation}
which are such that up to NNLO $M_8=M_4$, $M_9=M_5$, and
$M_{10}=M_7$. With these at hand, we find:
\begin{equation}
\footnotesize
\left[T\mathbb{K}^{(n_f)}T^{-1}\right]_{ij} =
\left\{
\begin{array}{lll}
M_1 & \quad i = 0 &\quad j = 0\\
\frac{1}{6}M_2 & \quad i = 0 &\quad  j = 1\\
0 & \quad i = 0 &\quad  2\leq j \leq n_f\\
-\frac{1}{n_f+1}M_3 & \quad i = 0 &\quad  j = n_f+1\\
\frac{1}{j(j-1)}M_2 & \quad i = 0 &\quad  n_f+ 2\leq j \leq 6\\
\\
M_4 & \quad i = 1&\quad  j=0\\
\frac16 M_5 & \quad i = 1 &\quad  j=1\\
0 & \quad i = 1 &\quad  2\leq j \leq n_f\\
-\frac{1}{n_f+1}M_6 & \quad i = 1 &\quad  j = n_f+1\\
\frac{1}{j(j-1)}M_5 & \quad i = 1 &\quad  n_f+
                                                            2\leq j
                                                            \leq 6\\
\\
\delta_{ij}M_7&\quad 2\leq i \leq n_f &\quad 0\leq j
                                                         \leq 6\\
\\
-n_fM_8 &\quad i = n_f+1 &\quad j =0\\
\frac{n_f}6\left(-M_9+(n_f+1)M_7\right) &\quad i = n_f+1 &\quad j =1\\
0 & \quad i = n_f+1 &\quad  2\leq j \leq n_f\\
\frac1{n_f+1}\left(n_f M_6+(n_f+1)M_7+M_5-M_9\right) &\quad i = n_f+1 &\quad j =n_f+1\\
\frac{n_f}{j(j-1)}\left( -M_9+(n_f+1)M_7\right) &\quad i = n_f+1
                  &\quad n_f+2 \leq j \leq 6\\
\\
M_4 & \quad n_f+2\leq i \leq 6 &\quad  j=0\\
\frac16 M_5 & \quad n_f+2\leq i \leq 6 &\quad  j=1\\
0 & \quad n_f+2\leq i \leq 6 &\quad  2\leq j \leq n_f\\
-\frac{1}{n_f+1}M_6 & \quad n_f+2\leq i \leq 6 &\quad  j = n_f+1\\
\frac{1}{j(j-1)}M_5 & \quad n_f+2\leq i \leq 6 &\quad  n_f+
                                                            2\leq j
                                                            \leq 6\\
\end{array}
\right.
\label{eq:algorithmEvN3LO}
\end{equation}

In the light of Eq.~(\ref{eq:IntMatchingCondsN3LO}), valence and
momentum sum rules imply the following identities for the matching
functions:
\begin{equation}
\begin{array}{l}
\displaystyle\int_0^1dy\,K_{ll}^{-,[n]}(y)=0\,,\\
\\
\displaystyle\int_0^1dy\,y\left[K_{gg}^{[n]}(y)+n_fK_{lg}^{[n]}(y)+K_{hg}^{[n]}(y)\right]=0\,,\\
\\
\displaystyle\int_0^1dy\,y\left[K_{gl}^{[n]}(y)+K_{ll}^{[n]}(y)+n_fK_{ll'}^{[n]}(y)+K_{hl}^{[n]}(y)\right]=0\,,\\
\\
\displaystyle\int_0^1dy\,y\left[K_{gh}^{[n]}(y)+K_{hh}^{[n]}(y)\right]=0\,,
\end{array}
\label{eq:QCDmatchingSumeRules}
\end{equation}
valid at any give order $n$ in the strong coupling $\alpha_s$.



\section{Evolution and matching in the physical basis}

In view of the implementation of QED corrections and the ensuing
complications due to the isospin symmetry breaking, we work out here
how to implement PDF evolution and matching in pure QCD but in the
physical basis. The purpose is to test the performance of this basis
in {\tt APFEL++} with the purpose of using this basis also in the
presence of QED corrections. Rather than the actual physical basis
$\{\overline{t}, \overline{b},\dots, \overline{d},g,d,\dots,b,t\}$, it
is convenient to use the basis
$\{f_i\}=\{t^-,b^-,\dots,d^-,g,d^+,\dots,b^+,t^+\}$, with
$q^{\pm} = q\pm\overline{q}$, in that it guarantees a partial
diagonalisation of the splitting-function matrix. In this basis, we
find that the evolution equations read:
\begin{equation}
\begin{array}{rcl}\displaystyle\frac{dq_k^-}{d\ln\mu^2} &=&
                                                            \displaystyle
                                                            P^{\rm
                                                            NV}q_k^-+P^{\rm
                                                            PV}\sum_{i=1,i\neq
                                                            k}^{n_f} q_i^-\,,\\
\\
\displaystyle\frac{dg}{d\ln\mu^2} &=&\displaystyle P_{gg} g+ P_{gq}
                                      \sum_{i=1}^{n_f} q_i^+\,,\\
\\
\displaystyle\frac{dq_k^+}{d\ln\mu^2} &=& \displaystyle 2P_{qg}
                                          g+P^{\rm
                                          NS}q_k^++P^{\rm
                                          PS}\sum_{i=1,i\neq k}^{n_f} q_i^+\,,
\end{array}
\end{equation}
where we have defined:
\begin{equation}
\begin{array}{l}
\displaystyle P^{\rm PV} = \frac{1}{n_f}(P^V-P^-)\,,\qquad P^{\rm PS} =
\frac{1}{n_f}(P_{qq}-P^+)\,,\\
\\
\displaystyle  P^{\rm NV} = P^-+P^{\rm PV}\,,\qquad P^{\rm NS} = P^++P^{\rm PS}\,.
\end{array}
\end{equation}
and with $k=1,\dots,n_f$ while:
\begin{equation}
\frac{dq_k^-}{d\ln\mu^2}=\frac{dq_k^+}{d\ln\mu^2} =
0\,,\quad\mbox{for}\quad k=n_f+1,\dots,6\,.
\end{equation}
Writing the full set of evolution equations as:
\begin{equation}
\frac{df_i}{d\ln\mu^2} = \mathbb{P}_{ij}^{(n_f)}f_j\,,
\end{equation}
with $i,j=-6,\dots,6$, one has that:
\begin{equation}
\mathbb{P}_{ij}^{(n_f)}=
\left\{
\begin{array}{lll}
P^{\rm PV} &\quad -n_f \leq i \leq -1 & \quad -n_f \leq j \leq -1\; ||\; j\neq i\\
P^{\rm NV}&\quad -n_f \leq i \leq -1 & \quad j=i\\
P_{gg} & \quad i = 0 &\quad j = 0\\
P_{gq} &\quad i = 0 & \quad 1 \leq j \leq n_f\\
2P_{qg} &\quad 1 \leq i \leq n_f & \quad j=0\\
P^{\rm NS}&\quad 1 \leq i \leq n_f & \quad j=i\\
P^{\rm PS} &\quad 1 \leq i \leq n_f & \quad 1 \leq j \leq n_f\; || \; j\neq i\\
0             & \quad\mbox{otherwise}
\end{array}
\right.\,.
\label{eq:splittingalgPhys}
\end{equation}

As far as matching conditions are concerned, using
Eq.~(\ref{eq:IntMatchingCondsN3LO}), one immediately finds:
\begin{equation}
\mathbb{K}^{(n_f)}_{ij} =
\left\{
\begin{array}{lll}
K_{gg} & \quad i = 0 &\quad j = 0\\
K_{gl} &\quad i = 0 & \quad 1 \leq j \leq n_f\\
K_{gh} &\quad i = 0 & \quad j = n_f + 1\\
K_{lg} &\quad 1 \leq i \leq n_f & \quad j=0\\
K_{ll}&\quad 1 \leq i \leq n_f & \quad j=i\\
K_{ll'} &\quad 1 \leq i \leq n_f & \quad 1 \leq j \leq n_f\\
K_{hg} &\quad i = n_f+1 & \quad j=0\\
K_{hl} &\quad i = n_f+1 & \quad 1\leq j\leq n_f\\
K_{hh} &\quad i = n_f+1 & \quad j=n_f+1\\
0             & \quad\mbox{elsewhere}
\end{array}
\right.\,.
\label{eq:algorithmPhys}
\end{equation}

It is useful to establish a connection with the $M$'s in
Eq.~(\ref{eq:McombsN3LO}). Specifically we find:
\begin{equation}
\begin{array}{l}
\displaystyle K_{gg} = M_1\,,\\
\\
\displaystyle K_{gl} =\frac{M_2-M_3}{n_f+1}\,,\\
\\
\displaystyle K_{gh} =\frac{M_2+n_fM_3}{n_f+1}\,,\\
\\
\displaystyle K_{lg} =\frac{M_4-M_8}{n_f+1}\,,\\
\\
\displaystyle K_{ll}=M_7\,,\\
\\
\displaystyle K_{ll'}=\frac{M_5-M_9}{n_f(n_f+1)}\,,\\
\\
\displaystyle K_{hg}=\frac{M_4+n_fM_8}{n_f+1}\,,\\
\\
\displaystyle K_{hl} =\frac{M_9-M_6}{n_f+1}-M_7\,,\\
\\
\displaystyle K_{hh}=\frac{M_5+n_fM_6}{n_f+1}\,,\\
\\
\displaystyle K_{ll}^-=M_{10}\,.
\end{array}
\end{equation}

\section{Evolution and matching with QED corrections}

In this section, we address the question of deriving the PDF evolution
operator in the presence of QED corrections including both photon and
charged-lepton PDFs. The fact that the photon coupling to quarks is
proportional to the electric charge, while the QCD coupling is
proportional to the strong coupling, complicates the description. The
general form of the QCD+QED evolution equations is:
\begin{equation}
\frac{\partial}{\partial t}
\begin{pmatrix}
L_\alpha\\
U_j \\
D_i\\
\Gamma_x\\
\overline{D}_k\\
\overline{U}_h\\
\overline{L}_\beta
\end{pmatrix} =
\begin{pmatrix}
\mathcal{P}_{L_\alpha L_\gamma} & \mathcal{P}_{L_\alpha U_e} & \mathcal{P}_{L_\alpha D_l}& \mathcal{P}_{L_\alpha \Gamma_y} & \mathcal{P}_{L_\alpha \overline{D}_m} & \mathcal{P}_{L_\alpha \overline{U}_n} & \mathcal{P}_{L_\alpha \overline{L}_\delta} \\ 
\mathcal{P}_{U_jL_\gamma} & \mathcal{P}_{U_jU_e} & \mathcal{P}_{U_jD_l} & \mathcal{P}_{U_j\Gamma_y} & \mathcal{P}_{U_j\overline{D}_m} & \mathcal{P}_{U_j\overline{U}_n} & \mathcal{P}_{U_j\overline{L}_\delta} \\ 
\mathcal{P}_{D_iL_\gamma} & \mathcal{P}_{D_iU_e} & \mathcal{P}_{D_iD_l} & \mathcal{P}_{D_i\Gamma_y} & \mathcal{P}_{D_i\overline{D}_m} & \mathcal{P}_{D_i\overline{U}_n} & \mathcal{P}_{D_i\overline{L}_\delta} \\ 
\mathcal{P}_{\Gamma_x L_\gamma} & \mathcal{P}_{\Gamma_x U_e} & \mathcal{P}_{\Gamma_x D_l} & \mathcal{P}_{\Gamma_x\Gamma_y} & \mathcal{P}_{\Gamma_x\overline{D}_m} & \mathcal{P}_{\Gamma_x\overline{U}_n} & \mathcal{P}_{\Gamma_x\overline{L}_\delta}\\
\mathcal{P}_{\overline{D}_kL_\gamma} & \mathcal{P}_{\overline{D}_kU_e} & \mathcal{P}_{\overline{D}_kD_l} & \mathcal{P}_{\overline{D}_k\Gamma_y} & \mathcal{P}_{\overline{D}_k\overline{D}_m} & \mathcal{P}_{\overline{D}_k\overline{U}_n} & \mathcal{P}_{\overline{D}_k\overline{L}_\delta}\\
\mathcal{P}_{\overline{U}_hL_\gamma} & \mathcal{P}_{\overline{U}_hU_e} & \mathcal{P}_{\overline{U}_hD_l} & \mathcal{P}_{\overline{U}_h\Gamma_y} & \mathcal{P}_{\overline{U}_h\overline{D}_m} & \mathcal{P}_{\overline{U}_h\overline{U}_n} & \mathcal{P}_{\overline{U}_h\overline{L}_\delta}\\
\mathcal{P}_{\overline{L}_\beta L_\gamma} & \mathcal{P}_{\overline{L}_\beta U_e} & \mathcal{P}_{\overline{L}_\beta D_l} & \mathcal{P}_{\overline{L}_\beta \Gamma_y} & \mathcal{P}_{\overline{L}_\beta \overline{D}_m} & \mathcal{P}_{\overline{L}_\beta \overline{U}_n} & \mathcal{P}_{\overline{L}_\beta \overline{L}_\delta}
\end{pmatrix}
\begin{pmatrix}
L_\gamma\\
U_e \\
D_l\\
\Gamma_y\\
\overline{D}_m\\
\overline{U}_n\\
\overline{L}_\delta
\end{pmatrix}\,,
\end{equation}
with $t=\ln\mu^2$ and where $D_i\in\{D_d,D_s,D_b\}$,
$U_j\in\{U_u,U_c,U_t\}$,
$L_\alpha\in \{L_{e^-},L_{\mu^-},L_{\tau^-}\}$ (and respective
anti-flavours such that $\overline{X}_i=X_{\overline{i}}$), and
$\Gamma_x\in\{g,\gamma\}$.  A sum over repeated indices is also
understood. Importantly, each distribution $X_i$ is in a given colour
state, such that $N_X$ copies of it exist, where $N_X$ is the number
of possible colour states. Clearly, we have that $N_D=N_U=N_c=3$ and
$N_L=1$. A decomposition of the splitting functions in terms of
valence and singlet contributions is as follows:
\begin{equation}
\begin{array}{l}
\mathcal{P}_{X_iY_j} = \mathcal{P}_{\overline{X}_i\overline{Y}_j} = \delta_{ij} \mathcal{P}_{XX}^{V}+\mathcal{P}_{XY}^{S}\,,\\
\mathcal{P}_{\overline{X}_iY_j} = \mathcal{P}_{X_i\overline{Y}_j}  = \delta_{ij} \mathcal{P}_{\overline{X}X}^{V}+\mathcal{P}_{\overline{X}Y}^{S}\,,\\
\mathcal{P}_{X_i\Gamma_x}  =\mathcal{P}_{\overline{X}_i\Gamma_x} =\mathcal{P}_{X\Gamma_x} \,,\\
\mathcal{P}_{\Gamma_x X_i}  =\mathcal{P}_{\Gamma_x\overline{X}_i} =\mathcal{P}_{\Gamma_x X} \,.
\end{array}
\end{equation}
where $X,Y\in\{D,U,L\}$. This decomposition implies that:
\begin{equation}
\begin{array}{rcl}
\displaystyle 
\frac{\partial}{\partial t}
\begin{pmatrix}
L_\alpha\\
U_j \\
D_i\\
\Gamma_x\\
\overline{D}_k\\
\overline{U}_h\\
\overline{L}_\beta
\end{pmatrix} &=&
\displaystyle 
\begin{pmatrix}
\mathcal{P}_{LL}^V & 0 & 0 & \mathcal{P}_{L\Gamma_y} & 0 & 0 & \mathcal{P}_{\overline{L}L}^V \\ 
0 & \mathcal{P}_{UU}^V & 0 & \mathcal{P}_{U\Gamma_y} & 0 & \mathcal{P}_{\overline{U}U}^V & 0 \\ 
0 & 0 & \mathcal{P}_{DD}^V & \mathcal{P}_{D\Gamma_y} & \mathcal{P}_{\overline{D}D}^V & 0 & 0 \\ 
\mathcal{P}_{\Gamma_x L} & \mathcal{P}_{\Gamma_x U} & \mathcal{P}_{\Gamma_x D} & \mathcal{P}_{\Gamma_x\Gamma_y} & \mathcal{P}_{\Gamma_xD} & \mathcal{P}_{\Gamma_xU} & \mathcal{P}_{\Gamma_xL}\\
0 & 0 & \mathcal{P}_{\overline{D}D}^V & \mathcal{P}_{D\Gamma_y} & \mathcal{P}_{DD}^V & 0 & 0\\
0 & \mathcal{P}_{\overline{U}U}^V & 0 & \mathcal{P}_{U\Gamma_y} & 0 & \mathcal{P}_{UU}^V & 0\\
\mathcal{P}_{\overline{L}L}^V & 0 & 0 & \mathcal{P}_{L\Gamma_y} & 0 & 0 & \mathcal{P}_{LL}^V
\end{pmatrix}
\begin{pmatrix}
L_\alpha\\
U_j \\
D_i\\
\Gamma_y\\
\overline{D}_k\\
\overline{U}_h\\
\overline{L}_\beta
\end{pmatrix}\\
\\
&+&
\displaystyle 
\begin{pmatrix}
\mathcal{P}^S_{LL} & \mathcal{P}^S_{UL} & \mathcal{P}^S_{DL}& 0 &\mathcal{P}^S_{\overline{D}L} & \mathcal{P}^S_{\overline{U}L} & \mathcal{P}^S_{\overline{L}L} \\ 
\mathcal{P}^S_{UL} & \mathcal{P}^S_{UU} & \mathcal{P}^S_{DU} & 0 & \mathcal{P}^S_{\overline{D}U} & \mathcal{P}^S_{\overline{U}U} & \mathcal{P}^S_{\overline{U}L} \\ 
\mathcal{P}^S_{DL} & \mathcal{P}^S_{DU} & \mathcal{P}^S_{DD} & 0 & \mathcal{P}^S_{\overline{D}D} & \mathcal{P}^S_{\overline{D}U} & \mathcal{P}^S_{\overline{D}L} \\ 
0 & 0 & 0 & 0 & 0 & 0 & 0\\
\mathcal{P}^S_{\overline{D}L} & \mathcal{P}^S_{\overline{D}U} & \mathcal{P}^S_{\overline{D}D} & 0 & \mathcal{P}^S_{DD} & \mathcal{P}^S_{DU} & \mathcal{P}^S_{DL}\\
\mathcal{P}^S_{\overline{U}L} & \mathcal{P}^S_{\overline{U}U} & \mathcal{P}^S_{\overline{D}U} & 0 & \mathcal{P}^S_{DU} & \mathcal{P}^S_{UU} & \mathcal{P}^S_{UL}\\
\mathcal{P}^S_{\overline{L}L} & \mathcal{P}^S_{\overline{U}L} &\mathcal{P}^S_{\overline{D}L} & 0 & \mathcal{P}^S_{DL} & \mathcal{P}^S_{UL} & \mathcal{P}^S_{LL}
\end{pmatrix}
\begin{pmatrix}
\sum_\gamma L_\gamma\\
\sum_e U_e \\
\sum_l D_l\\
\Gamma_y\\
\sum_m \overline{D}_m\\
\sum_n \overline{U}_n\\
\sum_\delta \overline{L}_\delta
\end{pmatrix}\,.
\end{array}
\end{equation}
Now, we define $X_i^\pm=X_i\pm\overline{X}_i$, which implies:
\begin{equation}
\small
\begin{array}{rcl}
\displaystyle 
\frac{\partial}{\partial t}
\begin{pmatrix}
L_\alpha^+\\
U_j^+ \\
D_i^+\\
\Gamma_x\\
D_i^-\\
U_j^-\\
L_\alpha^-
\end{pmatrix} &=&
\displaystyle 
\begin{pmatrix}
\mathcal{P}_{LL}^+ & 0 & 0 & 2\mathcal{P}_{L\Gamma_y} & 0 & 0 & 0 \\ 
0 & \mathcal{P}_{UU}^+ & 0 & 2\mathcal{P}_{U\Gamma_y} & 0 & 0 & 0 \\ 
0 & 0 & \mathcal{P}_{DD}^+ & 2\mathcal{P}_{D\Gamma_y} & 0 & 0 & 0 \\ 
\mathcal{P}_{\Gamma_x L} & \mathcal{P}_{\Gamma_x U} & \mathcal{P}_{\Gamma_x D} & \mathcal{P}_{\Gamma_x\Gamma_y} & 0 & 0 & 0\\
0 & 0 & 0 & 0 & \mathcal{P}_{DD}^- & 0 & 0\\
0 & 0 & 0 & 0 & 0 & \mathcal{P}_{UU}^- & 0\\
0 & 0 & 0 & 0 & 0 & 0 & \mathcal{P}_{LL}^-
\end{pmatrix}
\begin{pmatrix}
L_\alpha^+\\
U_j^+\\
D_i^+\\
\Gamma_y\\
D_i^-\\
U_j^-\\
L_\alpha^-
\end{pmatrix}\\
\\
&+&
\displaystyle 
\begin{pmatrix}
\mathcal{P}^{\rm PS}_{LL} & \mathcal{P}^{\rm PS}_{LU} & \mathcal{P}^{\rm PS}_{LD}& 0 &0 &0 & 0\\ 
\mathcal{P}^{\rm PS}_{UL} & \mathcal{P}^{\rm PS}_{UU} & \mathcal{P}^{\rm PS}_{UD} & 0 &0 &0 &0 \\ 
\mathcal{P}^{\rm PS}_{DL} & \mathcal{P}^{\rm PS}_{DU} & \mathcal{P}^{\rm PS}_{DD} & 0 & &0 &0 \\ 
0 & 0 & 0 & 0 & 0 & 0 & 0\\
0& 0&0 & 0 & \mathcal{P}^{\rm PV}_{DD} & \mathcal{P}^{\rm PV}_{DU} & \mathcal{P}^{\rm PV}_{DL}\\
0& 0& 0& 0 & \mathcal{P}^{\rm PV}_{UD} & \mathcal{P}^{\rm PV}_{UU} & \mathcal{P}^{\rm PV}_{UL}\\
0& 0&0 & 0 & \mathcal{P}^{\rm PV}_{LD} & \mathcal{P}^{\rm PV}_{LU} & \mathcal{P}^{\rm PV}_{LL}
\end{pmatrix}
\begin{pmatrix}
\sum_\gamma^{n_L} L_\gamma^+\\
\sum_e^{n_U} U_e^+ \\
\sum_l^{n_D} D_l^+\\
\Gamma_y\\
\sum_l^{n_D} D_l^-\\
\sum_e^{n_U} U_e^-\\
\sum_\gamma^{n_L} L_\gamma^-
\end{pmatrix}\,,
\end{array}
\end{equation}
where we have defined:
\begin{equation}
\mathcal{P}_{XX}^\pm=\mathcal{P}_{XX}^V\pm\mathcal{P}_{\overline{X}X}^V\,,\quad
\mathcal{P}_{XY}^{\rm PS} = \mathcal{P}^S_{XY}+\mathcal{P}^S_{\overline{X}Y}\,,\quad \mathcal{P}_{XY}^{\rm PV} = \mathcal{P}^S_{XY}-\mathcal{P}^S_{\overline{X}Y}\,.
\end{equation}
Also notice that the sums over flavours run up to $n_D$, $n_U$, and
$n_L$, which are the number of active down-type flavours, up-type
flavours, and leptons, respectively, and over all possible colour
states (despite not explicitly indicated).  We can now start
scrutinising the splitting functions. First, their perturbative
expansion reads:
\begin{equation}
\mathcal{P}= \underbrace{a_s
  \mathcal{P}^{[1,0]}+a\mathcal{P}^{[0,1]}}_{\mbox{LO}}+\underbrace{a_s^2\mathcal{P}^{[2,0]}+a_sa\mathcal{P}^{[1,1]}+a_s^2\mathcal{P}^{[0,2]}}_{\mbox{NLO}}+\underbrace{a_s^3\mathcal{P}^{[3,0]}+a_s^4\mathcal{P}^{[4,0]}}_{\mbox{NNLO
+ N$^3$LO (QCD only)}}\,,
\end{equation}
with $a_s=\alpha_s/(4\pi)$ and $a=\alpha/(4\pi)$ and where we have
grouped the perturbative terms according to the cumulative powers of
$a_s$ and $a$. Also, we are interested in QED corrections up to two
loops, \textit{i.e.} NLO. Higher orders are only accounted for in QCD
and are treated as discussed in the previous section. In this
approximation, the pure-valence block ($\mathcal{P}^{\rm PV}$), which
starts at NNLO, is purely given by QCD so that:
\begin{equation}
\mathcal{P}^{\rm PV}_{DL} = \mathcal{P}^{\rm PV}_{LD}=\mathcal{P}^{\rm PV}_{UL}=\mathcal{P}^{\rm PV}_{LU} =
\mathcal{P}^{\rm PV}_{LL} = 0\,,
\end{equation}
and:
\begin{equation}
 \mathcal{P}^{\rm PV}_{DD} = \mathcal{P}^{\rm PV}_{DU} =
 \mathcal{P}^{\rm PV}_{DD} = P^{\rm PV}\,.
\end{equation}
Moreover, at this order there is no coupling between leptons and
gluon, such that:
\begin{equation}
  \mathcal{P}_{Lg} = \mathcal{P}_{gL} = 0\,,
\end{equation}
so that, we finally have:
\begin{equation}
\small
\begin{array}{rcl}
\displaystyle 
\frac{\partial}{\partial t}
\begin{pmatrix}
L_\alpha^+\\
U_j^+ \\
D_i^+\\
g\\
\gamma\\
D_i^-\\
U_j^-\\
L_\alpha^-
\end{pmatrix} &=&
\displaystyle 
\begin{pmatrix}
\mathcal{P}_{LL}^+ & 0 & 0 & 0 & 2\mathcal{P}_{L\gamma} & 0 & 0 & 0 \\ 
0 & \mathcal{P}_{UU}^+ & 0 & 2\mathcal{P}_{Ug} & 2\mathcal{P}_{U\gamma} & 0 & 0 & 0 \\ 
0 & 0 & \mathcal{P}_{DD}^+ & 2\mathcal{P}_{Dg} & 2\mathcal{P}_{D\gamma} & 0 & 0 & 0 \\ 
0 & \mathcal{P}_{gU} & \mathcal{P}_{gD} & \mathcal{P}_{gg} & \mathcal{P}_{g\gamma}& 0 & 0 & 0\\
\mathcal{P}_{\gamma L} & \mathcal{P}_{\gamma U} & \mathcal{P}_{\gamma D} & \mathcal{P}_{\gamma g} & \mathcal{P}_{\gamma\gamma} & 0 & 0 & 0\\
0 & 0 & 0 & 0 & 0 & \mathcal{P}_{DD}^- & 0 & 0\\
0 & 0 & 0 & 0 & 0 & 0 & \mathcal{P}_{UU}^- & 0\\
0 & 0 & 0 & 0 & 0 & 0 & 0 & \mathcal{P}_{LL}^-
\end{pmatrix}
\begin{pmatrix}
L_\alpha^+\\
U_j^+\\
D_i^+\\
g\\
\gamma\\
D_i^-\\
U_j^-\\
L_\alpha^-
\end{pmatrix}\\
\\
&+&
\displaystyle 
\begin{pmatrix}
\mathcal{P}^{\rm PS}_{LL} & \mathcal{P}^{\rm PS}_{LU} & \mathcal{P}^{\rm PS}_{LD}& 0  & 0&0 &0 & 0\\ 
\mathcal{P}^{\rm PS}_{UL} & \mathcal{P}^{\rm PS}_{UU} & \mathcal{P}^{\rm PS}_{UD} & 0  & 0&0 &0 &0 \\ 
\mathcal{P}^{\rm PS}_{DL} & \mathcal{P}^{\rm PS}_{DU} & \mathcal{P}^{\rm PS}_{DD} & 0  & 0& &0 &0 \\ 
0 & 0 & 0 & 0 & 0 & 0 & 0 & 0\\
0 & 0 & 0 & 0 & 0 & 0 & 0 & 0\\
0& 0&0 & 0  & 0& P^{\rm PV} & P^{\rm PV} & 0\\
0& 0& 0& 0  & 0& P^{\rm PV} & P^{\rm PV} & 0\\
0& 0&0 & 0  & 0& 0 & 0 & 0
\end{pmatrix}
\begin{pmatrix}
\sum_\gamma^{n_L} L_\gamma^+\\
\sum_e^{n_U} U_e^+ \\
\sum_l^{n_D} D_l^+\\
g\\
\gamma\\
\sum_l^{n_D} D_l^-\\
\sum_e^{n_U} U_e^-\\
\sum_\gamma^{n_L} L_\gamma^-
\end{pmatrix}\,,
\end{array}
\end{equation}

It remains to assign to each non-zero entry of the splitting-function
matrices the correct value, which we will do below. In view of this
analysis, the basis that provides a partial diagonalisation of the
evolution equation along with a relative simple implementation is:
\begin{equation}
\{f_i\}=\{L_\tau^+, L_\mu^+,L_e^+,U_t^+,U_c^+,U_u^+,D_b^+,D_s^+,D_d^+,g,\gamma,D_d^-,D_s^-,D_b^-,U_u^-,U_c^-,U_t^-,L_e^-,L_\mu^-,L_\tau^-\}\,.
\label{eq:QCDQEDbasis}
\end{equation}
In principle, one could establish a basis that maximally diagonalises
the system (as for instance done in {\tt
  QCD\_QED\_common\_basis.pdf}). However, this basis is particularly
cumbersome to handle.

Let us now turn to threshold matching. In order to implement it, one
needs to be able establish what are the light flavours, what is the
heavy flavour, and what are the super-heavy flavours. To this purpose,
it is necessary to order distributions according to the thresholds of
the corresponding partons. Unfortunately, at variance with the pure
QCD case, this ordering is not reflected in the basis in
Eq.~(\ref{eq:QCDQEDbasis}). Given the ordered vector of thresholds
$\{\mu_i\}$, one can order the distributions accordingly:
\begin{equation}
  \{\widetilde{f}_i\}=\{X_i^+,\Gamma_x,X_i^-\}\,,
\label{eq:physbasismassordered}
\end{equation}
where a possible realistic ordering could be:
\begin{equation}
\{\widetilde{f}_i\}=\{U_t^+, D_b^+, L_\tau^+, U_c^+, L_\mu^+, D_s^+, D_d^+, U_u^+,L_e^+,g,\gamma, L_e^-,U_u^-, D_d^-,  D_s^-, L_\mu^-, U_c^-, L_\tau^-,D_b^-, U_t^-\}\,,
\end{equation}
which can however be determined dynamically according to the values of
thresholds. With this new basis, the matching at the $n_T$-th mass
threshold can be implemented sequentially, as in the case of pure
QCD. Using $l$ $(1\leq l \leq n_T)$, $h$ ($=n_T+1$), and $H$
($n_T+2\leq H \leq 9$) to indicate the light, heavy, and super-heavy
flavours, respectively, the matching is implemented through the
system:\footnote{It should be noted that, once the ordering in
  Eq.~(\ref{eq:physbasismassordered}) has been established on the
  basis of thresholds, the value of $n_T$ uniquely defines the values
  of $n_D$, $n_U$ and $n_L$ (which are in turn such that
  $n_D+n_U+n_L=n_T$).}
\begin{equation}
  \widetilde{f}_i^{(n_T+1)} = \left[\delta_{ij}+\mathcal{K}_{ij}(\kappa_h)\right] \widetilde{f}_j^{(n_T)}\,.
\end{equation}
with $\kappa_h=\mu_h/m_h$, being $m_h$ the mass (not the threshold) of
the heavy flavour $h$, and where:
\begin{equation}
\widetilde{f}_i^{(n_T)}=\{H^{+(n_T)},  h^{+(n_T)}, l^{+(n_T)}, \Gamma_x^{(n_T)}, l^{-(n_T)}, h^{-(n_T)}, H^{-(n_T)}\}\,,
\end{equation}
and similarly for $\widetilde{f}_i^{(n_T+1)}$. Given the
transformation matrix $T$, such that
$f_i=T_{ij}\widetilde{f}_j$,\footnote{Note that $T$ will be a sparse
  matrix with only one entry equal to one for each row and for each
  column but not along the diagonal.} one also has:
\begin{equation}
f_i^{(n_T+1)} = \left[\delta_{ij}+T_{ia}\mathcal{K}_{ab}(\kappa_h)T_{bj}^{-1}\right]f_j^{(n_T)}\,.
\end{equation}

One explicitly has:
\begin{equation}
\begin{array}{rcl}
  H^{+(n_T+1)} &=& H^{+(n_T)}\,,\\
\\
   h^{+(n_T+1)}&=&\displaystyle \sum_y\mathcal{K}_{hy}\Gamma_y^{(n_T)} + \sum_{l'}\mathcal{K}_{hl'}{l'}^{+(n_T)}+[1+\mathcal{K}_{hh}]h^{+(n_T)}\,,\\
  \\
   l^{+(n_T+1)}&=&\displaystyle\sum_y\mathcal{K}_{ly}\Gamma_y^{(n_T)}+[1+\mathcal{K}_{{\rm NS},l}]l^{+(n_T)}+\sum_{l'}\mathcal{K}_{ll'} {l'}^{+(n_T)}\,,\\
  \\
   \Gamma_x^{(n_T+1)}&=&\displaystyle\sum_y[\delta_{xy}+\mathcal{K}_{xy}]\Gamma_y^{(n_T)}+\sum_{l'}\mathcal{K}_{xl'} {l'}^{+(n_T)} + \mathcal{K}_{xh}h^{+(n_T)}\,,\\
  \\
  l^{-(n_T+1)} &=& [1+\mathcal{K}_{{\rm NS},l}^-]l^{-,(n_T)}\,,\\
\\
  h^{-(n_T+1)} &=& h^{-(n_T)}\,,\\
\\
  H^{-(n_T+1)} &=& H^{-(n_T)}\,,
\end{array}
\end{equation}
where, for brevity, we omitted the dependence of $\kappa_h$ of the
matching $\mathcal{K}$. Also, we notice that
$\mathcal{K}_{{\rm NS},l}^-$ differs from $\mathcal{K}_{{\rm NS},l}$
starting from $\mathcal{O}(\alpha_s^3)$.

As in the case of splitting functions, we will include QED corrections
up to NLO accuracy, leaving QCD only at higher orders. The
perturbative expansion of the matching matrix will thus be:
\begin{equation}
\mathcal{K}_{ij}=\underbrace{a_s \mathcal{K}_{ij}^{[1,0]}+a\mathcal{K}_{ij}^{[0,1]}}_{\mbox{NLO}}+\underbrace{a_s^2 \mathcal{K}_{ij}^{[2,0]}+a_s^3\mathcal{K}_{ij}^{[3,0]}}_{\mbox{NNLO + N$^3$LO (QCD only)}}\,.
\end{equation}
Under this assumption, the structure of matching greatly
simplifies. Indeed, up to NLO, which is the order up which QED
corrections are included, the only matching functions to be different
from zero are $\mathcal{K}_{xy}(\propto \delta_{xy})$,
$\mathcal{K}_{xh}$, $\mathcal{K}_{hy}$, and
$\mathcal{K}_{hh}$, all others start at NNLO or beyond. Since
corrections beyond NLO are QCD only, they do not carry any flavour
dependence and any function that involves leptons and/or photon is
identically zero. This leads to:
\begin{equation}
\begin{array}{rcl}
  H^{+(n_T+1)} &=& H^{+(n_T)}\,,\\
\\
   h^{+(n_T+1)}&=&\displaystyle \mathcal{K}_{hg}g^{(n_T)} + \mathcal{K}_{h \gamma}\gamma^{(n_T)} +\mathcal{K}_{hl'}\sum_{l'} {l'}^{+(n_T)}+[1+\mathcal{K}_{hh}]h^{+(n_T)}\,,\\
  \\
   l^{+(n_T+1)}&=&\displaystyle \mathcal{K}_{lg}g^{(n_T)}+[1+\mathcal{K}_{{\rm NS},l}]l^{+(n_T)}+\mathcal{K}_{ll'} \sum_{l'} {l'}^{+(n_T)}\,,\\
  \\
   g^{(n_T+1)}&=&\displaystyle [1+\mathcal{K}_{gg}]g^{(n_T)}+\mathcal{K}_{gl'} \sum_{l'}{l'}^{+(n_T)} + \mathcal{K}_{gh}h^{+(n_T)}\,,\\
\\
   \gamma^{(n_T+1)}&=&\displaystyle[1+\mathcal{K}_{\gamma\gamma}]\gamma^{(n_T)}+\mathcal{K}_{\gamma h}h^{+(n_T)}\,,\\
  \\
  l^{-(n_T+1)} &=& [1+\mathcal{K}_{{\rm NS},l}^-]l^{-,(n_T)}\,,\\
\\
  h^{-(n_T+1)} &=& h^{-(n_T)}\,,\\
\\
  H^{-(n_T+1)} &=& H^{-(n_T)}\,.
\end{array}
\end{equation}

If $h$ is a lepton, such that $h=L_\alpha$, we have that all
distributions of the basis $f_i$ match on themselves
($f_i^{(n_T+1)}=f_i^{(n_T)}$), except:
\begin{equation}
\begin{array}{rcl}
 \displaystyle
  L_\alpha^{+(n_T+1)}&=&\displaystyle[1+\mathcal{K}_{LL}]L_\alpha^{+(n_T)}+\mathcal{K}_{L\gamma}\gamma^{(n_T)}\,,\\
\\
  \gamma^{(n_T+1)}&=&\displaystyle \mathcal{K}_{\gamma L} L_\alpha^{+(n_T)} +[1+\mathcal{K}_{\gamma\gamma}^L]\gamma^{(n_T)}\,,
\end{array}
\label{eq:MatchingLeptons}
\end{equation}
where all functions $\mathcal{K}$ implicitly depend on the variable
$\kappa_{L_\alpha}$. If instead $h=X_i=D_i,U_i$, we would still have
$f_i^{(n_T+1)}=f_i^{(n_T)}$ for all but:
\begin{equation}
\begin{array}{rcl}
   X_i^{+(n_T+1)}&=&\displaystyle \mathcal{K}_{Qg}g^{(n_T)} + \mathcal{K}_{X\gamma}\gamma^{(n_T)} +\mathcal{K}_{Qq'}\left(\sum_{j=1}^{n_D} D_j^{+(n_T)}+\sum_{j=1}^{n_U} U_j^{+(n_T)}\right)+[1+\mathcal{K}_{XX}]X_i^{+(n_T)}\,,\\
  \\
   q^{+(n_T+1)}&=&\displaystyle \mathcal{K}_{qg}g^{(n_T)}+[1+\mathcal{K}_{{\rm NS},q}]q^{+(n_T)}+\mathcal{K}_{qq'}\left(\sum_{j=1}^{n_D} D_j^{+(n_T)}+\sum_{j=1}^{n_U} U_j^{+(n_T)}\right)\,,\\
  \\
   g^{(n_T+1)}&=&\displaystyle [1+\mathcal{K}_{gg}]g^{(n_T)}+\mathcal{K}_{gq}\left(\sum_{j=1}^{n_D} D_j^{+(n_T)}+\sum_{j=1}^{n_U} U_j^{+(n_T)}\right) + \mathcal{K}_{gQ}X_i^{+(n_T)}\,,\\
\\
   \gamma^{(n_T+1)}&=&\displaystyle[1+\mathcal{K}_{\gamma\gamma}^X]\gamma^{(n_T)}+\mathcal{K}_{\gamma X}X_i^{+(n_T)}\,,\\
  \\
  q^{-(n_T+1)} &=& [1+\mathcal{K}_{{\rm NS},q}^-]q^{-,(n_T)}\,,
\end{array}
\label{eq:MatchingQuarks}
\end{equation}
where all functions $\mathcal{K}$ implicitly depend on the variable
$\kappa_{X_i}$ and where we used the character $Q$ ($q$) to mean a
generic heavy (light) quark, no matter whether of up or down type.

In fact, Eqs.~(\ref{eq:MatchingLeptons}) and~(\ref{eq:MatchingQuarks})
can be put together by allowing the index $X$ in
Eq.~(\ref{eq:MatchingQuarks}) to run over all possible values
($X=D,U,L$) and by setting:
\begin{equation}
  \mathcal{K}_{Qg}=\mathcal{K}_{Qq'}=\mathcal{K}_{qg}=\mathcal{K}_{{\rm
      NS},q}=\mathcal{K}_{qq'}=\mathcal{K}_{gg}=\mathcal{K}_{gq}=\mathcal{K}_{gQ}=0\,,
\label{eq:IfLeptonMatching}
\end{equation}
when $h=L_\alpha$.

We are now in a position to give the specific form of splitting and
matching functions. Let us start with splitting functions. We have 30
splitting functions to define, they are:
\begin{equation}
\begin{array}{l}
\mathcal{P}_{DD}^+,\mathcal{P}_{UU}^+,\mathcal{P}_{LL}^+,\\
\mathcal{P}_{DD}^-,\mathcal{P}_{UU}^-,\mathcal{P}_{LL}^-,\\
\mathcal{P}^{\rm PS}_{LL}, \mathcal{P}^{\rm PS}_{LU}, \mathcal{P}^{\rm PS}_{LD}\mathcal{P}^{\rm PS}_{UL},\mathcal{P}^{\rm PS}_{UU},\mathcal{P}^{\rm PS}_{UD},\mathcal{P}^{\rm PS}_{DL},\mathcal{P}^{\rm PS}_{DU},\mathcal{P}^{\rm PS}_{DD},\\
P^{\rm PV},\\
\mathcal{P}_{Dg},\mathcal{P}_{Ug},\mathcal{P}_{D\gamma},\mathcal{P}_{U\gamma},\mathcal{P}_{L\gamma},\\
\mathcal{P}_{gD},\mathcal{P}_{gU},\mathcal{P}_{\gamma D},\mathcal{P}_{\gamma U},\mathcal{P}_{\gamma L},\\
\mathcal{P}_{gg},\mathcal{P}_{g\gamma},\mathcal{P}_{\gamma g},\mathcal{P}_{\gamma\gamma}\,.
\end{array}
\end{equation}
At $\mathcal{O}(\alpha_s)$, we have that:
\begin{equation}
\begin{array}{l}
\mathcal{P}_{LL}^{+,[1,0]}=\mathcal{P}_{LL}^{-,[1,0]}=\mathcal{P}^{\rm PS,[1,0]}_{LL}=\mathcal{P}^{\rm PS,[1,0]}_{LU}=\mathcal{P}^{\rm PS,[1,0]}_{LD}=\mathcal{P}^{\rm
  PS,[1,0]}_{UL}=\mathcal{P}^{\rm PS,[1,0]}_{UU}=\mathcal{P}^{\rm
  PS,[1,0]}_{UD}=\mathcal{P}^{\rm PS,[1,0]}_{DL}=\mathcal{P}^{\rm PS,[1,0]}_{DD}=\\
\\
  \mathcal{P}^{\rm PS,[1,0]}_{DU}= P^{\rm PV [1,0]}=\mathcal{P}_{D\gamma}^{[1,0]}=\mathcal{P}_{U\gamma}^{[1,0]}=\mathcal{P}_{L\gamma}^{[1,0]}=\mathcal{P}_{\gamma
  D}^{[1,0]}=\mathcal{P}_{\gamma U}^{[1,0]}=\mathcal{P}_{\gamma
  L}^{[1,0]}=\mathcal{P}_{g\gamma}^{[1,0]}=\mathcal{P}_{\gamma
  g}^{[1,0]}=\mathcal{P}_{\gamma\gamma}^{[1,0]}=0\,,
\end{array}
\end{equation}
and:
\begin{equation}
\begin{array}{l}
\displaystyle
  \mathcal{P}_{DD}^{+,[1,0]}=\mathcal{P}_{UU}^{+,[1,0]}=\mathcal{P}_{DD}^{-,[1,0]}=\mathcal{P}_{UU}^{-,[1,0]}
  = 2C_F\left(\frac{1+x^2}{1-x}\right)_+\,,\\
\\
\displaystyle \mathcal{P}_{Dg}^{[1,0]}=\mathcal{P}_{Ug}^{[1,0]}=2T_R\left[x^2+(1-x)^2\right]\,,\\
\\
\displaystyle \mathcal{P}_{gD}^{[1,0]}=\mathcal{P}_{gU}^{[1,0]}=2C_F\frac{1+(1-x)^2}{x}\,,\\
\\
\displaystyle \mathcal{P}_{gg}^{[1,0]}=4C_A\left[\frac{x}{(1-x)_+}+\frac{1-x}{x}+x(1-x)\right]+\beta_0\delta(1-x)\,.
\end{array}
\end{equation}
At $\mathcal{O}(\alpha)$, we have an analogous picture:
\begin{equation}
\begin{array}{l}
\mathcal{P}^{\rm PS,[0,1]}_{LL}=\mathcal{P}^{\rm
  PS,[0,1]}_{LU}=\mathcal{P}^{\rm PS,[0,1]}_{LD}=\mathcal{P}^{\rm
  PS,[0,1]}_{UL}=\mathcal{P}^{\rm PS,[0,1]}_{UU}=\mathcal{P}^{\rm
  PS,[0,1]}_{UD}=\mathcal{P}^{\rm
  PS,[0,1]}_{DL}=\mathcal{P}^{\rm PS,[0,1]}_{DU}=\mathcal{P}^{\rm
  PS,[0,1]}_{DD}=\\
\\
  P^{\rm PV [0,1]}=\mathcal{P}_{Dg}^{[0,1]}=\mathcal{P}_{Ug}^{[0,1]}=\mathcal{P}_{gD}^{[0,1]}=\mathcal{P}_{gU}^{[0,1]}=\mathcal{P}_{gg}^{[0,1]}=\mathcal{P}_{g\gamma}^{[0,1]}=\mathcal{P}_{\gamma g}^{[0,1]}=0\,,
\end{array}
\end{equation}
and:
\begin{equation}
\begin{array}{l}
\displaystyle
  \mathcal{P}_{XX}^{+,[0,1]}=\mathcal{P}_{XX}^{-,[0,1]} = 2e_X^2\left(\frac{1+x^2}{1-x}\right)_+\,,\\
\\
\displaystyle \mathcal{P}_{X\gamma}^{[0,1]}=2e_X^2N_X\left[x^2+(1-x)^2\right]\,,\\
\\
\displaystyle \mathcal{P}_{\gamma X}^{[0,1]}=2e_X^2\frac{1+(1-x)^2}{x}\,,\\
\\
\displaystyle \mathcal{P}_{\gamma\gamma}^{[0,1]}=-\frac{4}{3}e_\Sigma^2\delta(1-x)\,,
\end{array}
\label{eq:OneLoopQEDSplittings}
\end{equation}
for $X=D,U,L$, with the charges being $e_D^2=1/9$, $e_U^2=4/9$,
$e_L^2=1$, and where we have defined:\footnote{Note that in
  Eq.~(\ref{eq:OneLoopQEDSplittings}) we have that $\alpha=2$, but
  below we will also need $\alpha=4$.}
\begin{equation}
  e_\Sigma^\alpha=\sum_{X=D,U,L}n_XN_Xe_X^\alpha\,.
\end{equation}

Now let us move to two loops considering first
$\mathcal{O}(\alpha_s^2)$. We will refrain from writing the explicit
expressions for the splitting functions and refer to the usual QCD
ones. At this order we have:
\begin{equation}
\begin{array}{l}
\mathcal{P}_{LL}^{+,[2,0]}=\mathcal{P}_{LL}^{-,[2,0]}=\mathcal{P}^{\rm PS,[2,0]}_{LL}=\mathcal{P}^{\rm PS,[2,0]}_{LU}=\mathcal{P}^{\rm PS,[2,0]}_{LD}=\mathcal{P}^{\rm PS,[2,0]}_{UL}=\mathcal{P}^{\rm PS,[2,0]}_{DL}=  P^{\rm PV [2,0]}=\\
\\
\mathcal{P}_{D\gamma}^{[2,0]}=\mathcal{P}_{U\gamma}^{[2,0]}=\mathcal{P}_{L\gamma}^{[2,0]}=\mathcal{P}_{\gamma
  D}^{[2,0]}=\mathcal{P}_{\gamma U}^{[2,0]}=\mathcal{P}_{\gamma
  L}^{[2,0]}=\mathcal{P}_{g\gamma}^{[2,0]}=\mathcal{P}_{\gamma
  g}^{[2,0]}=\mathcal{P}_{\gamma\gamma}^{[2,0]}=0\,,
\end{array}
\end{equation}
and:
\begin{equation}
\begin{array}{l}
\displaystyle
  \mathcal{P}_{DD}^{+,[2,0]}=\mathcal{P}_{UU}^{+,[2,0]}= P^{+,[1]}\,,\\
\\
\mathcal{P}_{DD}^{-,[2,0]}=\mathcal{P}_{UU}^{-,[2,0]} = P^{-,[1]}\,,\\
\\
\mathcal{P}^{\rm PS,[2,0]}_{DD}=\mathcal{P}^{\rm PS,[2,0]}_{DU}=\mathcal{P}^{\rm  PS,[2,0]}_{DU}=\mathcal{P}^{\rm  PS,[2,0]}_{UU}= P^{\rm PS,[1]}\,,\\
\\
\displaystyle \mathcal{P}_{Dg}^{[2,0]}=\mathcal{P}_{Ug}^{[2,0]}=P_{qg}^{[1]}\,,\\
\\
\displaystyle \mathcal{P}_{gD}^{[2,0]}=\mathcal{P}_{gU}^{[2,0]}=P_{gq}^{[1]}\,,\\
\\
\displaystyle \mathcal{P}_{gg}^{[2,0]}=P_{gg}^{[1]}\,.
\end{array}
\end{equation}

For the splitting functions at $\mathcal{O}(\alpha^2)$, we have:
\begin{equation}
\begin{array}{l}
  P^{\rm PV [0,2]}=\mathcal{P}_{Dg}^{[0,2]}=\mathcal{P}_{Ug}^{[0,2]}=\mathcal{P}_{gD}^{[0,2]}=\mathcal{P}_{gU}^{[0,2]}=\mathcal{P}_{gg}^{[0,2]}=\mathcal{P}_{g\gamma}^{[0,2]}=\mathcal{P}_{\gamma g}^{[0,2]}=0\,,
\end{array}
\end{equation}
and:
\begin{equation}
\begin{array}{rcl}
  \displaystyle
  \mathcal{P}_{XX}^{\pm,[0,2]}&=&\displaystyle e_X^4\bigg\{-4\bigg[\left(2\ln(x)\ln(1-x)+\frac{3}{2}\ln(x)\right)\frac{1+x^2}{1-x}+\frac{3+7x}{2}\ln(x)\\
\\
&+&\displaystyle\frac{1+x}{2}\ln^2(x)+5(1-x)+\left(\frac{\pi^2}{2}-\frac{3}{8}-6\zeta_3\right)\delta(1-x)\bigg]\\
\\
&-&\displaystyle \frac{e_\Sigma^2}{e_X^2}4\left[\frac{4}{3}(1-x)+\frac{1+x^2}{(1-x)_+}\left(\frac{2}{3}\ln(x)+\frac{10}{9}\right)+\left(\frac{2\pi^2}{9}+\frac{1}{6}\right)\delta(1-x)\right]\\
\\
&\pm&\displaystyle 4\left[4(1-x)+2(1+x)\ln(x)+2\frac{1+x^2}{1+x}S_2(x)\right]\bigg\}\,,\\
\\
  \displaystyle \mathcal{P}^{\rm PS,[0,2]}_{XY}&=&\displaystyle N_Xe_X^2e_Y^28\left[\frac{20}{9x}-2+6x-\frac{56}{9}x^2+\left(1+5x+\frac{8}{3}x^2\right)\ln(x)-(1+x)\ln^2(x)\right]\,,\\
  \\
  \displaystyle \mathcal{P}_{X\gamma}^{[0,2]}&=&\displaystyle N_X e_X^4
  2\bigg[4-9x-(1-4x)\ln(x)-(1-2x)\ln^2(x)+4\ln(1-x)\\
\\
&+&\displaystyle (x^2+(1-x)^2)\left(2\ln^2\left(\frac{1-x}{x}\right)-4\ln\left(\frac{1-x}{x}\right)-\frac{2\pi^2}{3}+10 \right)\bigg]\,,\\
  \\
  \displaystyle \mathcal{P}_{\gamma X}^{[0,2]}&=&\displaystyle
                                                  e_X^4\bigg\{4\bigg[-(3\ln(1-x)+\ln^2(1-x))\frac{1+(1-x)^2}{x}+\left(2+\frac{7}{2}x\right)\ln(x)-\left(1-\frac{x}{2}\right)\ln^2(x)\\
\\
&-&\displaystyle 2x\ln(1-x)-\frac{7}{2}x-\frac{5}{2}\bigg]- \frac{e_\Sigma^2}{e_X^2}4\left[\frac{4}{3}x+\frac{1+(1-x)^2}{x}\left(\frac{20}{9}+\frac{4}{3}\ln(1-x)\right)\right]\bigg\}\,,\\
  \\
  \displaystyle \mathcal{P}_{\gamma\gamma}^{[0,2]}&=&\displaystyle e_\Sigma^44\left[-16+8x+\frac{20}{3}x^2+\frac{4}{3x}-(6+10x)\ln(x)-2(1+x)\ln^2(x)-\delta(1-x)\right]\,,
\end{array}
\end{equation}
The expressions are taken from Ref.~\cite{deFlorian:2016gvk}.

Let us now move to the $\mathcal{O}(\alpha_s\alpha)$ splitting
functions, whose relevant expressions can be read off from
Ref.~\cite{deFlorian:2015ujt}. At this order, we find:
\begin{equation}
\begin{array}{l}
\mathcal{P}_{LL}^{+,[1,1]}=\mathcal{P}_{LL}^{-,[1,1]}=\mathcal{P}^{\rm PS,[1,1]}_{LL}=\mathcal{P}^{\rm
  PS,[1,1]}_{LU}=\mathcal{P}^{\rm PS,[1,1]}_{LD}=\mathcal{P}^{\rm
  PS,[1,1]}_{UL}=\mathcal{P}^{\rm PS,[1,1]}_{UU}=\\
  \\
  \mathcal{P}^{\rm
  PS,[1,1]}_{UD}=\mathcal{P}^{\rm
  PS,[1,1]}_{DL}=\mathcal{P}^{\rm PS,[1,1]}_{DU}=\mathcal{P}^{\rm
  PS,[1,1]}_{DD}= P^{\rm PV [1,1]}=\mathcal{P}_{L\gamma}^{[1,1]}=\mathcal{P}_{\gamma L}^{[1,1]}=0\,,
\end{array}
\end{equation}
and:
\begin{equation}
  \begin{array}{rcl}
    \displaystyle \mathcal{P}_{XX}^{\pm,[1,1]}&=&\displaystyle
                                                  -e_X^28C_F\bigg[\left(2\ln(1-x)+\frac{3}{2}\right)\ln(x)\frac{1+x^2}{1-x}+\frac{3+7x}{2}\ln(x)+\frac{1+x}{2}\ln^2(x)\\
    \\
    &+&\displaystyle 5(1-x)+\left(\frac{\pi^2}{2}-\frac{3}{8}-6\zeta_3\right)\delta(1-x)\bigg]\\
    \\
    &\pm&\displaystyle
          e_X^28C_F\left[4(1-x)+2(1+x)\ln(x)+\frac{1+x^2}{1+x}S_2(x)\right]\,,
          \quad X=U,D\\
    \\
    \displaystyle \mathcal{P}_{gX}^{[1,1]}&=&\displaystyle
                                              e_X^24C_F\bigg[-(3\ln(1-x)+\ln^2(1-x))\frac{1+(1-x)^2}{x}+\left(2+\frac{7}{2}x\right)\ln(x)\\
    \\
    &-&\displaystyle \left(1-\frac{x}{2}\right)\ln^2(x)-2x\ln(1-x)-\frac{7}{2}x-\frac{5}{2}\bigg]\,,\quad X=U,D\\
    \\
        \displaystyle \mathcal{P}_{\gamma X}^{[1,1]}&=&\displaystyle\mathcal{P}_{gX}^{[1,1]}\,,\quad X=U,D\\
    \\

    \displaystyle \mathcal{P}_{Xg}^{[1,1]}&=&\displaystyle
                                              e_X^22T_R\bigg[4-9x-(1-4x)\ln(x)-(1-2x)\ln^2(x)+4\ln(1-x)\\
    \\
    &+&\displaystyle (x^2+(1-x)^2)\left(2\ln^2\left(\frac{1-x}{x}\right)-4\ln\left(\frac{1-x}{x}\right)-\frac{2\pi^2}{3}+10\right)\bigg]\,,\quad X=U,D\\
    \\
        \displaystyle \mathcal{P}_{X\gamma}^{[1,1]}&=&\displaystyle e_X^22N_XC_F\bigg[4-9x-(1-4x)\ln(x)-(1-2x)\ln^2(x)+4\ln(1-x)\\
    \\
    &+&\displaystyle (x^2+(1-x)^2)\left(2\ln^2\left(\frac{1-x}{x}\right)-4\ln\left(\frac{1-x}{x}\right)-\frac{2\pi^2}{3}+10\right)\bigg]\,,\quad X=U,D\\
    \\
  \displaystyle \mathcal{P}_{gg}^{[1,1]}&=& \displaystyle-4T_R\left(n_De_D^2+n_Ue_U^2\right)\delta(1-x)\,,\\
\\
  \displaystyle \mathcal{P}_{g\gamma}^{[1,1]}&=& \displaystyle 4N_cC_F\left(n_De_D^2+n_Ue_U^2\right)\left[-16+8x+\frac{20}{3}x^2+\frac{4}{3x}-(6+10x)\ln(x)-2(1+x)\ln^2(x)\right]\,,\\
\\
  \displaystyle \mathcal{P}_{\gamma g}^{[1,1]}&=&\displaystyle 4T_R\left(n_De_D^2+n_Ue_U^2\right) \left[-16+8x+\frac{20}{3}x^2+\frac{4}{3x}-(6+10x)\ln(x)-2(1+x)\ln^2(x)\right]\,,\\
\\
  \displaystyle \mathcal{P}_{\gamma\gamma}^{[1,1]}&=&\displaystyle-4N_cC_F\left(n_De_D^2+n_Ue_U^2\right)\delta(1-x) \,.
\end{array}
\end{equation}

At higher orders, where only QCD corrections are included, one has
that all splitting functions involving leptons and/or photon are zero:
\begin{equation}
\begin{array}{l}
  \mathcal{P}_{LL}^{+,[n,0]}=\mathcal{P}_{LL}^{-,[n,0]}=\mathcal{P}^{{\rm
  PS} ,[n,0]}_{LL}= \mathcal{P}^{{\rm PS}
  ,[n,0]}_{LU}=\mathcal{P}^{{\rm PS} ,[n,0]}_{LD}=\mathcal{P}^{{\rm
  PS} ,[n,0]}_{UL}=\mathcal{P}^{{\rm PS}
  ,[n,0]}_{UL}=\mathcal{P}^{{\rm PS} ,[n,0]}_{DL}=\\
  \\
  \mathcal{P}_{D\gamma}^{[n,0]}=\mathcal{P}_{U\gamma}^{[n,0]}=\mathcal{P}_{L\gamma}^{[n,0]}=\mathcal{P}_{\gamma D}^{[n,0]}=\mathcal{P}_{\gamma U}^{[n,0]}=\mathcal{P}_{\gamma L}^{[n,0]}=\mathcal{P}_{g\gamma}^{[n,0]}=\mathcal{P}_{\gamma g}^{[n,0]}=\mathcal{P}_{\gamma\gamma}^{[n,0]}=0\,,
\end{array}
\end{equation}
for $n\geq3$. The others are given by the QCD ones as follows:
\begin{equation}
  \begin{array}{l}
    \mathcal{P}_{DD}^{\pm,[n,0]}=\mathcal{P}_{UU}^{\pm,[n,0]}=P^{\pm,[n-1]}\,,\\
\mathcal{P}^{{\rm PS},[n,0]}_{UU}=\mathcal{P}^{{\rm
    PS},[n,0]}_{UD}=\mathcal{P}^{{\rm
    PS},[n,0]}_{DU}=\mathcal{P}^{{\rm PS},[n,0]}_{DD}=P^{{\rm PS},[n-1]}\,,\\
P^{{\rm PV},[n,0]}=P^{{\rm PV},[n-1]}\,,\\
\mathcal{P}_{Dg}^{[n,0]}=\mathcal{P}_{Ug}^{[n,0]}=P_{qg}^{[n-1]}\,,\\
\mathcal{P}_{gD}^{[n,0]}=\mathcal{P}_{gU}^{[n,0]}=P_{gq}^{[n-1]}\,,\\
\mathcal{P}_{gg}^{[n,0]}=P_{gg}^{[n-1]}\,,
  \end{array}
\end{equation}
where $P^{[n-1]}$ denote the QCD splitting functions at N$^{n-1}$LO.

We can finally move to the matching conditions. Looking at
Eq.~(\ref{eq:MatchingQuarks}), staring fron NLO accuracy, we need
specify the various functions $\mathcal{K}$. At each heavy-quark
threshold, the functions $\mathcal{K}$ also depend on the ration
$\kappa_X=\mu_X/m_X$, where $\mu_X$ and $m_X$ are the threshold and
the mass of the quarks that is becoming active. We have a total of 12
functions to determine:
\begin{equation}
\mathcal{K}_{Qg},\mathcal{K}_{X\gamma},\mathcal{K}_{Qq'},\mathcal{K}_{XX},\mathcal{K}_{qg},\mathcal{K}_{{\rm NS},q},\mathcal{K}_{qq'},\mathcal{K}_{gg},\mathcal{K}_{gq},\mathcal{K}_{gQ},\mathcal{K}_{\gamma\gamma}^X,\mathcal{K}_{\gamma X}\,,
\end{equation}
As already discussed above, if $X$ is a lepton ($X=L$), many of these
vanish (see Eq.~(\ref{eq:IfLeptonMatching})). If instead $X$ is either
a down- or a up-type quark ($X=D,U$), the functions coincide with the
pure-QCD ones. Specifically, referring to
Eq.~(\ref{eq:IntMatchingCondsN3LO}), one has:
\begin{equation}
\begin{array}{rcl}
\displaystyle\mathcal{K}_{Qg}&=&\displaystyle \left\{
\begin{array}{ll}
K_{hg}&\quad X=D,U\,,\\
0 & \quad X=L\,,
\end{array}
\right.\\
\\
\displaystyle\mathcal{K}_{Qq'}&=&\displaystyle \left\{ 
\begin{array}{ll}
K_{hl}&\quad X=D,U\,,\\
0 & \quad X=L\,,
\end{array}
\right.\\
\\
\displaystyle\mathcal{K}_{qg}&=&\displaystyle \left\{ 
\begin{array}{ll}
K_{lg}&\quad X=D,U\,,\\
0 & \quad X=L\,,
\end{array}
\right.\\
\\
\displaystyle\mathcal{K}_{{\rm NS},q}&=&\displaystyle \left\{ 
\begin{array}{ll}
K_{ll}&\quad X=D,U\,,\\
0 & \quad X=L\,,
\end{array}
\right.\\
\\
\displaystyle\mathcal{K}_{qq'}&=&\displaystyle \left\{ 
\begin{array}{ll}
K_{ll'}&\quad X=D,U\,,\\
0 & \quad X=L\,,
\end{array}
\right.\\
\\
\displaystyle\mathcal{K}_{gg}&=&\displaystyle \left\{ 
\begin{array}{ll}
K_{gg}&\quad X=D,U\,,\\
0 & \quad X=L\,,
\end{array}
\right.\\
\\
\displaystyle\mathcal{K}_{gq}&=&\displaystyle \left\{ 
\begin{array}{ll}
K_{gl}&\quad X=D,U\,,\\
0 & \quad X=L\,,
\end{array}
\right.\\
\\
\displaystyle\mathcal{K}_{gQ}&=&\displaystyle \left\{ 
\begin{array}{ll}
K_{gh}&\quad X=D,U\,,\\
0 & \quad X=L\,,
\end{array}
\right.\\
\\
\displaystyle\mathcal{K}_{{\rm NS},q}^-&=&\displaystyle \left\{ 
\begin{array}{ll}
K_{ll}^-&\quad X=D,U\,,\\
0 & \quad X=L\,.
\end{array}
\right.\\
\\
\end{array}
\end{equation}

The remaining are:
\begin{equation}
\begin{array}{rcl}
\displaystyle \mathcal{K}_{XX} &=&\displaystyle
\left\{
\begin{array}{ll}
 K_{hh}&\quad X=D,U\,,\\
\\
\displaystyle 2e_X^2\left[\frac{1+x^2}{1-x}\left(\ln\left(\kappa_X^2\right)-1-2\ln(1-x)\right)\right]_+& \quad X=L\,,
\end{array}
\right.\\
\\
\displaystyle\mathcal{K}_{X\gamma}&=&\displaystyle 4N_Xe_X^2[x^2+(1-x)^2]\ln\left(\kappa_X^2\right)\,,\\
\\
\displaystyle\mathcal{K}_{\gamma X}&=&\displaystyle 2e_X^2\left[\frac{1+(1-x)^2}{x}\right]\left[\ln\left(\kappa_X^2\right)-1-2\ln(x)\right]\,,\\
\\
\displaystyle\mathcal{K}_{\gamma\gamma}^X &=&\displaystyle -\frac{4}{3}e_X^2 \ln\left(\kappa_X^2\right)\delta(1-x)\,.
\end{array}
\end{equation}

\subsection{Sum rules}

One of the possible check the implementation of the correctness of
splitting functions is that sum rules have to be fulfilled at all
scales order by order in perturbation theory. This check is
particularly useful when including QED corrections given the
proliferation of splitting functions. Sum rules can be checked
numerically either at the level of evolved PDFs (as long as they
fulfil them at the initial scale) or at the level of splitting
functions.

Let us start with the momentum sum rule, which at the level of PDFs,
prescribes that:
\begin{equation}
  \int_0^1dy\,y\left[\sum_{i}D_i^+(y,\mu)+\sum_{j}U_j^+(y,\mu)+\sum_{\alpha}L_\alpha^+(y,\mu)+g(y,\mu)+\gamma(y,\mu)\right]=1\,,
\end{equation}
for all values of $\mu$. Taking the derivative w.r.t. $\mu$ of this
equality and using the evolution equations allows us to derive
integral equations which must obeyed by splitting functions order by
order in perturbation theory. Specifically, we find:
\begin{equation}
\begin{array}{l}
\displaystyle
  \int_0^1dy\,y\left[\mathcal{P}_{LL}^{+,[a,b]}(y)+n_L\mathcal{P}_{LL}^{{\rm PS},[a,b]}(y) +n_U\mathcal{P}_{UL}^{{\rm PS},[a,b]}(y) +n_D\mathcal{P}_{DL}^{{\rm PS},[a,b]}(y)+\mathcal{P}_{\gamma L}^{[a,b]}(y)\right]=0\,,\\
\\
\displaystyle \int_0^1dy\,y\left[\mathcal{P}_{UU}^{+,[a,b]}(y) +n_L\mathcal{P}_{LU}^{{\rm PS},[a,b]}(y) +n_U\mathcal{P}_{UU}^{{\rm PS},[a,b]}(y) +n_D\mathcal{P}_{DU}^{{\rm PS},[a,b]}(y)+\mathcal{P}_{gU}^{[a,b]}(y)+\mathcal{P}_{\gamma U}^{[a,b]}(y)\right]=0\,,\\
\\
\displaystyle \int_0^1dy\,y\left[\mathcal{P}_{DD}^{+,[a,b]}(y) +n_L\mathcal{P}_{LD}^{{\rm PS},[a,b]}(y) +n_U\mathcal{P}_{UD}^{{\rm PS},[a,b]}(y) +n_D\mathcal{P}_{DD}^{{\rm PS},[a,b]}(y)+\mathcal{P}_{gD}^{[a,b]}(y)+\mathcal{P}_{\gamma D}^{[a,b]}(y)\right]=0\,,\\
\\
\displaystyle\int_0^1dy\,y\left[2n_U\mathcal{P}_{Ug}^{[a,b]}(y)+2n_D\mathcal{P}_{Dg}^{[a,b]}(y)+\mathcal{P}_{\gamma g}^{[a,b]}(y) +\mathcal{P}_{gg}^{[a,b]}(y)\right]=0\,,\\
\\
\displaystyle \int_0^1dy\,y\left[2n_L\mathcal{P}_{L\gamma}^{[a,b]}(y)+2n_U\mathcal{P}_{U\gamma}^{[a,b]}(y)+2n_D\mathcal{P}_{D\gamma}^{[a,b]}(y) +\mathcal{P}_{g\gamma}^{[a,b]}(y)+\mathcal{P}_{\gamma\gamma}^{[a,b]}(y)\right]=0\,,
\end{array}
\end{equation}
for the any of the perturbative orders $[a,b]$ of our interest. For
the leading-order sets of splitting functions, it is straightforward
to check analytically that this equalities are true. Beyond these
orders analytic checks are still possible but more cumbersome.

We next consider the valence (or counting) sum rule, which at the
level pf PDFs reads:
\begin{equation}
\int_0^1dy\,X_i^-(y,\mu)=\mbox{constant}\,,\quad X=D,U,L\,,
\end{equation}
where the constant depends on $X$. At the level of splitting
functions, the valence sum rule implies:
\begin{equation}
\begin{array}{l}
\displaystyle \int_0^1dy\,\left[\mathcal{P}_{XX}^{-,[a,b]}(y)+n_X P^{{\rm
  PV},[a,b]}(y)\right]=0\,,\quad X=D,U\,,\\
\\
\displaystyle \int_0^1dy\,\mathcal{P}_{LL}^{-,[a,b]}(y)=0\,.
\end{array}
\end{equation}

Also matching functions are subject to integral identities deriving
from the sum rules. Indeed, matching functions must be such that
valence and momentum sum rules are preserved when a threshold is
crossed. In pure QCD, we found the identities in
Eq.~(\ref{eq:QCDmatchingSumeRules}). In the presence of QED
corrections, we refer to the system in Eq.~(\ref{eq:MatchingQuarks})
and from the valence sum rule we derive:
\begin{equation}
\int_0^1dy\, \mathcal{K}_{{\rm NS},q}^{-,[a,b]}(y)=0\,,
\end{equation}
while from the momentum sum rule we find:
\begin{equation}
\begin{array}{l}
\displaystyle \int_0^1dy\,y\left[\mathcal{K}_{Qg}^{[a,b]}(y)+n_f\mathcal{K}_{qg}^{[a,b]}(y)+\mathcal{K}_{gg}^{[a,b]}(y)\right]=0\,,\\
\\
\displaystyle \int_0^1dy\,y\left[\mathcal{K}_{X\gamma}^{[a,b]}(y)+\mathcal{K}_{\gamma\gamma}^{X,[a,b]}(y)\right]=0\,,\\
\\
\displaystyle \int_0^1dy\,y\left[\mathcal{K}_{Qq'}^{[a,b]}(y)+\mathcal{K}_{{\rm NS},q}^{[a,b]}(y)+n_f\mathcal{K}_{qq'}^{[a,b]}(y)+\mathcal{K}_{gq}^{[a,b]}(y)\right]=0\,,\\
\\
\displaystyle \int_0^1dy\,y\left[\mathcal{K}_{XX}^{[a,b]}(y)+\mathcal{K}_{gQ}^{[a,b]}(y)+\mathcal{K}_{\gamma X}^{[a,b]}(y)\right]=0\,,
\end{array}
\end{equation}
where we defined $n_f=n_D+n_U$.


\newpage
\bibliographystyle{ieeetr}
\bibliography{bibliography}

\end{document}
